\documentclass[11pt,a4paper]{article}
\usepackage[T1]{fontenc}
\usepackage{amsmath,amsthm}
\usepackage{newtxtext,newtxmath}
\usepackage[margin=2.3cm,headheight=14pt]{geometry}
\usepackage{longtable,microtype,graphicx,booktabs,array,float,caption,subcaption,enumitem,xcolor,tikz,pgfplots,pgfplotstable,siunitx,fancyhdr,titlesec}
\usepackage[ruled,vlined,linesnumbered]{algorithm2e}
\usepackage[colorlinks=true,linkcolor=accent,citecolor=accent,urlcolor=accent]{hyperref}
\usepackage{aliascnt}
\usepackage[nameinlink,capitalise]{cleveref}
\emergencystretch=2em
% block paragraphs: no first-line indent, a half line between paragraphs, so every left edge lines up
\usepackage[skip=0.55\baselineskip plus 2pt minus 1pt]{parskip}
% let floats share pages with text instead of leaving half-empty pages
\renewcommand{\topfraction}{0.85}\renewcommand{\bottomfraction}{0.6}\renewcommand{\textfraction}{0.12}
\renewcommand{\floatpagefraction}{0.8}\setcounter{topnumber}{3}\setcounter{totalnumber}{4}
\setlength{\textfloatsep}{14pt plus 4pt minus 4pt}\setlength{\floatsep}{12pt plus 4pt minus 2pt}
\usepackage{needspace}
% pages holding only floats stack them from the top instead of spreading them out
\makeatletter\setlength{\@fptop}{0pt}\setlength{\@fpsep}{16pt plus 4pt}\setlength{\@fpbot}{0pt plus 1fil}\makeatother
\pgfplotsset{compat=1.18}
\usetikzlibrary{arrows.meta,positioning,calc,shapes}
\definecolor{accent}{HTML}{1F4E79}
\definecolor{soft}{HTML}{EAF1F8}
\definecolor{cG1}{HTML}{D55E00}\definecolor{cG2}{HTML}{009E73}\definecolor{cG3}{HTML}{0072B2}
\definecolor{cDP}{HTML}{CC79A7}\definecolor{cBB}{HTML}{555555}
\captionsetup{font=small,labelfont=bf,skip=6pt}
\setlist{itemsep=1pt,topsep=3pt}
\sisetup{group-separator={,},group-minimum-digits=4}
\titleformat{\section}{\Large\bfseries\color{accent}}{\thesection}{0.8em}{}
\titleformat{\subsection}{\large\bfseries}{\thesubsection}{0.7em}{}
\pagestyle{fancy}\fancyhf{}
\fancyhead[L]{\small Steiner trees on Fibonacci heaps}\fancyhead[R]{\small CS23I1006 \& CS23I1002}
\fancyfoot[C]{\small\thepage}
\renewcommand{\headrulewidth}{0.4pt}
% aliascnt gives each environment its own counter name, so cleveref prints Lemma/Corollary, not Theorem
% parskip zeroes the list \topsep that amsthm uses, so the spacing around theorems is set here
\newtheoremstyle{plain}{9pt}{5pt}{\itshape}{}{\bfseries}{.}{.5em}{}
\newtheoremstyle{definition}{9pt}{5pt}{\normalfont}{}{\bfseries}{.}{.5em}{}
\theoremstyle{plain}\newtheorem{theorem}{Theorem}[section]
\newaliascnt{lemma}{theorem}\newtheorem{lemma}[lemma]{Lemma}\aliascntresetthe{lemma}
\newaliascnt{corollary}{theorem}\newtheorem{corollary}[corollary]{Corollary}\aliascntresetthe{corollary}
\theoremstyle{definition}\newaliascnt{definition}{theorem}\newtheorem{definition}[definition]{Definition}\aliascntresetthe{definition}
\crefname{lemma}{Lemma}{Lemmas}\crefname{corollary}{Corollary}{Corollaries}\crefname{definition}{Definition}{Definitions}
\SetKwProg{Fn}{function}{}{}\SetKwComment{tcp}{\(\triangleright\) }{}
\SetAlFnt{\small}\SetAlCapFnt{\small}\SetAlCapNameFnt{\small}
\newcommand{\OPT}{\mathrm{OPT}}
\newcommand{\Dn}{D(n)}
% numbers come from results/*.csv via scripts/make_tables.py
\makeatletter
\newcommand{\val}[1]{\ifcsname v@#1\endcsname\csname v@#1\endcsname\else\PackageError{report}{missing value #1}{}\fi}
\makeatother
% generated by scripts/make_tables.py, do not edit
\expandafter\def\csname v@abl.DP.dense.ratio\endcsname{1.51}
\expandafter\def\csname v@abl.DP.sparse.ratio\endcsname{1.73}
\expandafter\def\csname v@abl.G1.dense.ratio\endcsname{1.10}
\expandafter\def\csname v@abl.G1.sparse.ratio\endcsname{2.08}
\expandafter\def\csname v@abl.G2.dense.ratio\endcsname{1.49}
\expandafter\def\csname v@abl.G2.sparse.ratio\endcsname{2.12}
\expandafter\def\csname v@abl.G3.dense.ratio\endcsname{1.75}
\expandafter\def\csname v@abl.G3.sparse.ratio\endcsname{1.89}
\expandafter\def\csname v@abl.dp.big.bin\endcsname{2.87}
\expandafter\def\csname v@abl.dp.big.decr\endcsname{2{,}263{,}578}
\expandafter\def\csname v@abl.dp.big.extr\endcsname{3{,}276{,}736}
\expandafter\def\csname v@abl.dp.big.fib\endcsname{3.85}
\expandafter\def\csname v@abl.dp.big.inst\endcsname{grid_k16_s0}
\expandafter\def\csname v@abl.dp.c12.bin\endcsname{77.9}
\expandafter\def\csname v@abl.dp.c12.decr\endcsname{259{,}028}
\expandafter\def\csname v@abl.dp.c12.extr\endcsname{255{,}529}
\expandafter\def\csname v@abl.dp.c12.fib\endcsname{169.3}
\expandafter\def\csname v@abl.dp.c16.bin\endcsname{2.0}
\expandafter\def\csname v@abl.dp.c16.decr\endcsname{6{,}319}
\expandafter\def\csname v@abl.dp.c16.extr\endcsname{7{,}511}
\expandafter\def\csname v@abl.dp.c16.fib\endcsname{7.5}
\expandafter\def\csname v@abl.dp.c17.bin\endcsname{151.3}
\expandafter\def\csname v@abl.dp.c17.decr\endcsname{184{,}598}
\expandafter\def\csname v@abl.dp.c17.extr\endcsname{255{,}518}
\expandafter\def\csname v@abl.dp.c17.fib\endcsname{223.8}
\expandafter\def\csname v@adv.g1.cost\endcsname{42}
\expandafter\def\csname v@adv.g1.opt\endcsname{30}
\expandafter\def\csname v@adv.g1.ratio\endcsname{1.400}
\expandafter\def\csname v@adv.g2.cost\endcsname{68}
\expandafter\def\csname v@adv.g2.opt\endcsname{49}
\expandafter\def\csname v@adv.g2.ratio\endcsname{1.388}
\expandafter\def\csname v@adv.g3.cost\endcsname{58}
\expandafter\def\csname v@adv.g3.opt\endcsname{48}
\expandafter\def\csname v@adv.g3.ratio\endcsname{1.208}
\expandafter\def\csname v@amort.ops\endcsname{7.43\,billion}
\expandafter\def\csname v@amort.ratio\endcsname{0.32}
\expandafter\def\csname v@amort.test.ratio\endcsname{0.25}
\expandafter\def\csname v@bb.solved\endcsname{152}
\expandafter\def\csname v@bb.solved.C\endcsname{19}
\expandafter\def\csname v@bb.solved.all\endcsname{198}
\expandafter\def\csname v@bb.tl\endcsname{300}
\expandafter\def\csname v@bb.total\endcsname{153}
\expandafter\def\csname v@bb.total.all\endcsname{199}
\expandafter\def\csname v@bb.unsolved\endcsname{1}
\expandafter\def\csname v@both.bbfaster\endcsname{97}
\expandafter\def\csname v@both.count\endcsname{173}
\expandafter\def\csname v@dp.attempted\endcsname{174}
\expandafter\def\csname v@dp.cfit\endcsname{9.926e-10}
\expandafter\def\csname v@dp.grid.k20\endcsname{115.4}
\expandafter\def\csname v@dp.maxtime\endcsname{159.3}
\expandafter\def\csname v@dp.mem.k17\endcsname{78.6}
\expandafter\def\csname v@dp.mem.k20\endcsname{503}
\expandafter\def\csname v@dp.solved\endcsname{173}
\expandafter\def\csname v@env.cores\endcsname{16}
\expandafter\def\csname v@env.cpu\endcsname{Intel Core Ultra 7 255H}
\expandafter\def\csname v@env.gxx\endcsname{g++  16.2.1 20261029 (Red Hat 16.2.1-3)}
\expandafter\def\csname v@env.mem\endcsname{31}
\expandafter\def\csname v@env.os\endcsname{Linux 7.2.14}
\expandafter\def\csname v@exact.checked\endcsname{279}
\expandafter\def\csname v@g1.maxgap\endcsname{74.17}
\expandafter\def\csname v@g1.meangap\endcsname{14.19}
\expandafter\def\csname v@g1.optcount\endcsname{18}
\expandafter\def\csname v@g2.maxgap\endcsname{74.17}
\expandafter\def\csname v@g2.meangap\endcsname{8.93}
\expandafter\def\csname v@g2.optcount\endcsname{63}
\expandafter\def\csname v@g3.b18.centres\endcsname{23}
\expandafter\def\csname v@g3.b18.scanned\endcsname{176{,}400}
\expandafter\def\csname v@g3.maxgap\endcsname{20.83}
\expandafter\def\csname v@g3.meangap\endcsname{4.65}
\expandafter\def\csname v@g3.optcount\endcsname{60}
\expandafter\def\csname v@gap.Adversarial.G1\endcsname{50.48}
\expandafter\def\csname v@gap.Adversarial.G2\endcsname{48.78}
\expandafter\def\csname v@gap.Adversarial.G3\endcsname{1.60}
\expandafter\def\csname v@gap.B.G1\endcsname{2.38}
\expandafter\def\csname v@gap.B.G2\endcsname{0.18}
\expandafter\def\csname v@gap.B.G3\endcsname{0.65}
\expandafter\def\csname v@gap.C.G1\endcsname{5.36}
\expandafter\def\csname v@gap.C.G2\endcsname{1.99}
\expandafter\def\csname v@gap.C.G3\endcsname{2.33}
\expandafter\def\csname v@gap.Grid.G1\endcsname{2.89}
\expandafter\def\csname v@gap.Grid.G2\endcsname{0.54}
\expandafter\def\csname v@gap.Grid.G3\endcsname{0.57}
\expandafter\def\csname v@gap.I080.G1\endcsname{18.51}
\expandafter\def\csname v@gap.I080.G2\endcsname{10.73}
\expandafter\def\csname v@gap.I080.G3\endcsname{8.14}
\expandafter\def\csname v@gap.P6E.G1\endcsname{2.77}
\expandafter\def\csname v@gap.P6E.G2\endcsname{0.36}
\expandafter\def\csname v@gap.P6E.G3\endcsname{0.09}
\expandafter\def\csname v@gap.Random.G1\endcsname{7.95}
\expandafter\def\csname v@gap.Random.G2\endcsname{1.03}
\expandafter\def\csname v@gap.Random.G3\endcsname{2.42}
\expandafter\def\csname v@gap.count\endcsname{199}
\expandafter\def\csname v@gap.max.steinlib\endcsname{41.0}
\expandafter\def\csname v@inst.generated\endcsname{46}
\expandafter\def\csname v@inst.steinlib\endcsname{153}
\expandafter\def\csname v@inst.total\endcsname{199}
\expandafter\def\csname v@rows.total\endcsname{995}
\expandafter\def\csname v@sph.worse_than_g1\endcsname{48}
\expandafter\def\csname v@time.C.G1\endcsname{0.03}
\expandafter\def\csname v@time.C.G2\endcsname{1.70}
\expandafter\def\csname v@time.C.G3\endcsname{0.15}


\begin{document}
% ------------------------------------------------------------------ cover
\begin{titlepage}
\begin{tikzpicture}[remember picture,overlay]
\fill[accent] (current page.north west) rectangle ([yshift=-7.2cm]current page.north east);
\node[anchor=north west,text=white,align=left,font=\sffamily] at ([xshift=2.3cm,yshift=-1.8cm]current page.north west)
{\large Advanced Data Structures and Algorithms \(\cdot\) Course Assignment\\[18pt]
 {\fontsize{30}{34}\selectfont\bfseries The Steiner Tree Problem}\\[10pt]
 {\Large on Fibonacci heaps: greedy, dynamic programming, branch-and-bound}};
\end{tikzpicture}
\vspace*{7.2cm}

\noindent\begin{tabular}{@{}p{4.2cm}p{10.5cm}@{}}
\textbf{Team} & Sri Sooraj S \hfill CS23I1006\\
 & Nishaant Bhaskar \hfill CS23I1002\\[6pt]
\textbf{Institute} & IIITDM Kancheepuram\\[6pt]
\textbf{Problem} & \#7 Steiner Tree Problem (optimisation version)\\[6pt]
\textbf{Data structure} & Fibonacci heap (Fredman--Tarjan), with full amortized analysis\\[6pt]
\textbf{Greedy strategies} & G1 distance-network MST (Kou--Markowsky--Berman)\\
 & G2 shortest-path heuristic (Takahashi--Matsuyama)\\
 & G3 best-gain triples (Zelikovsky, ratio \(11/6\))\\[6pt]
\textbf{Exact paradigms} & Dynamic programming (Dreyfus--Wagner with a heap-based grow step)\\
 & Branch-and-bound (reductions, dual ascent, reduced-cost elimination)\\[6pt]
\textbf{Instances} & SteinLib B, C, P6E, I080 (\val{inst.steinlib} instances with published optima),\\
 & \val{inst.generated} generated instances (adversarial and scaling)\\[6pt]
\textbf{Code} & C++17, standard library only; \texttt{CS23I1006\_CS23I1002.zip}
\end{tabular}

\vfill
\noindent\colorbox{soft}{\parbox{\dimexpr\linewidth-2\fboxsep}{\small
\textbf{Results in one paragraph.}
Both exact methods agree with every published SteinLib optimum they reach.
Branch-and-bound proves optimality on \val{bb.solved} of \val{bb.total} SteinLib instances
(all of B, P6E and I080, \val{bb.solved.C} of 20 in C) within \val{bb.tl}\,s each; the dynamic program solves
\val{dp.solved} of the \val{dp.attempted} instances with \(k\le 20\).
Of the greedy strategies, Zelikovsky's triple rule (G3) has the smallest mean gap to the optimum over all
\val{gap.count} instances (\val{g3.meangap}\%, against \val{g1.meangap}\% for G1 and \val{g2.meangap}\% for G2), while G2,
which tries every root, is best on four of the seven sets at about ten times the cost of G3.
Over \val{amort.ops} heap operations inside the solvers, the measured work is
\val{amort.ratio} of the amortized bound proved in \cref{sec:amortized}.
}}
\end{titlepage}

\tableofcontents
\newpage

% ------------------------------------------------------------------ 1
\section{The problem}\label{sec:problem}

\begin{definition}[Steiner tree problem, optimisation version]
An \emph{instance} is a triple \(I=(G,w,T)\) where \(G=(V,E)\) is a connected undirected graph with
\(n=|V|\) and \(m=|E|\), \(w:E\to\mathbb{Z}_{\ge0}\) is an edge weight, and \(T\subseteq V\) is the set of
\(k=|T|\) \emph{terminals}. The instance space is
\[\mathcal I=\{(G,w,T): G \text{ connected},\ w\ge 0,\ \emptyset\neq T\subseteq V\}.\]
The \emph{feasible set} \(F(I)\) is the set of trees \(S=(V_S,E_S)\) with \(S\subseteq G\) and \(T\subseteq V_S\).
The \emph{objective} is \(f:F(I)\to\mathbb{R}_{\ge0}\), \(f(S)=\sum_{e\in E_S}w(e)\), and the task is to find
\(S^\ast\in\arg\min_{S\in F(I)}f(S)\). We write \(\OPT(I)=f(S^\ast)\). Vertices of \(V_S\setminus T\) are
\emph{Steiner vertices}.
\end{definition}

\noindent The objective maps into \(\mathbb{R}^+\) as the assignment asks: every feasible tree has a non-negative cost, and the cost is
positive unless all terminals coincide or are joined by zero-weight edges.
Two special cases are polynomial: \(k=2\) is a shortest path and \(T=V\) is a minimum spanning tree. In general the
decision version is NP-complete (Karp, 1972), and it stays NP-complete on planar graphs and on graphs with unit
weights. It is also APX-hard: no polynomial algorithm reaches a ratio below \(96/95\) unless P\,=\,NP
(Chleb\'ik and Chleb\'ikov\'a, 2008). So the report studies approximation (\cref{sec:greedy}) and two exact
algorithms whose running time is exponential only in a parameter: \(k\) for the dynamic program (\cref{sec:dp}) and
the number of branching decisions for branch-and-bound (\cref{sec:bnb}).

Throughout, \(d(u,v)\) is the shortest-path distance in \(G\) and \(D=(T,\binom{T}{2},d)\) is the \emph{distance
network}: the complete graph on the terminals weighted by \(d\). Every algorithm below works on \(D\) or on
shortest paths in \(G\), and that is where the data structure enters.

\paragraph{Where the Fibonacci heap is used.}
Every one of the five solvers spends most of its time in Dijkstra or Prim runs, and both are sequences of
\textsc{Insert}, \textsc{Decrease-Key} and \textsc{Extract-Min}. A Dijkstra run on a graph with \(n\) vertices and
\(m\) edges does \(n\) inserts, \(n\) extracts and up to \(m\) decrease-keys. With a binary heap every operation costs
\(O(\log n)\) and the run costs \(O(m\log n)\); with a Fibonacci heap \textsc{Decrease-Key} is \(O(1)\) amortized and the
run costs \(O(m+n\log n)\). \Cref{tab:usage} lists the calls. The branch-and-bound uses the heap twice more: as
the queue of open search nodes, keyed by lower bound, and as the queue of terminals waiting for a dual ascent step.

\begin{table}[H]\centering\small
\begin{tabular}{@{}l>{\raggedright\arraybackslash}p{4.9cm}l>{\footnotesize}l@{}}\toprule
Solver & Heap calls per run & Total time & File (\texttt{src/})\\\midrule
G1 KMB & \(k\) Dijkstra + Prim on \(D\) + Prim on a subgraph & \(O(k(m+n\log n))\) & \texttt{greedy1\_kmb.cpp}\\
G2 SPH & \(k-1\) Dijkstra per root, \(k\) roots & \(O(k^2(m+n\log n))\) & \texttt{greedy2\_sph.cpp}\\
G3 Zelikovsky & \(k\) Dijkstra + one Prim on \(D\) per round & \(O(k(m+n\log n)+k^4+k^3n)\) & \texttt{greedy3\_zelikovsky.cpp}\\
DP & one Dijkstra per subset of \(k-1\) terminals & \(O(3^{k}n+2^{k}(m+n\log n))\) & \texttt{dp\_dreyfus\_wagner.cpp}\\
B\&B & node queue, dual-ascent queue, two reduced-cost Dijkstra runs per node & exponential (\cref{sec:bnb}) & \texttt{bnb\_steiner.cpp}\\\bottomrule
\end{tabular}
\caption{Where each solver calls the Fibonacci heap (\texttt{include/fib\_heap.hpp}).}\label{tab:usage}
\end{table}

% ------------------------------------------------------------------ 2
\section{The data structure: a Fibonacci heap}\label{sec:heap}

\subsection{Why this structure}
The state that all five algorithms manage is a set of tentative labels (a distance, a key in Prim, a lower bound
in branch-and-bound) from which the smallest is removed again and again, while other labels only go down. The
number of label decreases is much larger than the number of removals: Dijkstra removes each vertex once but may
lower a label once per edge. In the dynamic program of \cref{sec:dp} the grow step runs Dijkstra from
\(n\) sources at once on every one of \(2^{k-1}\) layers, so a single solve performs millions of
\textsc{Decrease-Key} calls (\cref{tab:heapcounts}). The Fibonacci heap is built for exactly this ratio: it makes
\textsc{Decrease-Key} \(O(1)\) amortized by postponing all restructuring to \textsc{Extract-Min}.

\begin{table}[H]\centering\small
\begin{tabular}{@{}lccccc@{}}\toprule
Operation & Binary heap & Binomial heap & Pairing heap & \multicolumn{2}{c}{Fibonacci heap}\\
 & worst & worst & amortized & worst & amortized\\\midrule
\textsc{Make-Heap}, \textsc{Minimum} & \(O(1)\) & \(O(1)\) & \(O(1)\) & \(O(1)\) & \(O(1)\)\\
\textsc{Insert} & \(O(\log n)\) & \(O(\log n)\) & \(O(1)\) & \(O(1)\) & \(O(1)\)\\
\textsc{Union} & \(O(n)\) & \(O(\log n)\) & \(O(1)\) & \(O(1)\) & \(O(1)\)\\
\textsc{Extract-Min} & \(O(\log n)\) & \(O(\log n)\) & \(O(\log n)\) & \(O(n)\) & \(O(\log n)\)\\
\textsc{Decrease-Key} & \(O(\log n)\) & \(O(\log n)\) & \(o(\log n)\)\textsuperscript{a} & \(O(n)\)\textsuperscript{b} & \(O(1)\)\\
\textsc{Delete} & \(O(\log n)\) & \(O(\log n)\) & \(O(\log n)\) & \(O(n)\) & \(O(\log n)\)\\\midrule
Dijkstra, \(n\) vertices, \(m\) edges & \(O(m\log n)\) & \(O(m\log n)\) & \(o(m\log n)\) & & \(O(m+n\log n)\)\\\bottomrule
\end{tabular}
\caption{Priority queues compared. \textsuperscript{a}The exact amortized bound for pairing heaps is open;
\(O(1)\) is impossible in the standard model (Fredman, 1999). \textsuperscript{b}A cascade of cuts can follow a path
as long as the tallest tree, which can reach \(\Theta(n)\).}\label{tab:pq}
\end{table}

The Fibonacci heap is the only structure in the table that gives Dijkstra and Prim the
\(O(m+n\log n)\) bound, and that bound is what the complexity results of \cref{sec:greedy,sec:dp} rest on.
\Cref{sec:ablation} measures whether the asymptotic advantage turns into wall-clock time on our instances.

\subsection{Representation}
A Fibonacci heap is a collection of heap-ordered rooted trees. Each node \(x\) stores
\texttt{key}, \texttt{item} (the payload: a vertex or a search-node id), \texttt{parent}, \texttt{child} (any one
child), \texttt{left}/\texttt{right} (siblings in a circular doubly linked list), \texttt{degree} (number of
children) and a boolean \texttt{mark}. The roots form the \emph{root list}, another circular list, and the heap keeps a
pointer \texttt{min} to the root with the smallest key, together with \(n\) (number of nodes), \(t(H)\) (number of roots) and
\(m(H)\) (number of marked nodes). A node is marked when it has lost a child since it last became the child of another node.

In \texttt{include/fib\_heap.hpp} nodes come from a chunked pool, so a handle is a stable pointer and
\textsc{Clear} recycles the memory without freeing it; the dynamic program calls \textsc{Clear} once per subset layer
instead of allocating \(n\) new nodes. The class \texttt{FibPQ} wraps the heap for graph algorithms: it keeps one
handle per vertex so that \textsc{Decrease-Key}\((v,k)\) can be called with a vertex id. The heap also counts its
own work (links, cuts, cascading cuts, roots scanned, children promoted), which \cref{sec:amortized-check} uses to
check the amortized analysis against real runs.

\subsection{Abstract data type}
\Cref{tab:adt} lists the operations. Fundamental operations are those of a mergeable min-priority queue;
auxiliary operations are internal and are listed because the amortized analysis charges their cost separately.

\begin{table}[H]\centering\small
\begin{tabular}{@{}p{3.9cm}p{6.1cm}cc@{}}\toprule
Operation & Effect & Worst case & Amortized\\\midrule
\multicolumn{4}{@{}l}{\emph{Fundamental}}\\
\textsc{Make-Heap}() & empty heap & \(O(1)\) & \(O(1)\)\\
\textsc{Insert}\((H,x)\) & add node \(x\) with key \(x.\mathit{key}\) & \(O(1)\) & \(O(1)\)\\
\textsc{Minimum}\((H)\) & node with the smallest key & \(O(1)\) & \(O(1)\)\\
\textsc{Extract-Min}\((H)\) & remove and return the minimum & \(O(n)\) & \(O(\log n)\)\\
\textsc{Union}\((H_1,H_2)\) & meld two heaps, \(H_2\) becomes empty & \(O(1)\) & \(O(1)\)\\
\textsc{Decrease-Key}\((H,x,k)\) & set \(x.\mathit{key}\gets k\le x.\mathit{key}\) & \(O(n)\) & \(O(1)\)\\
\textsc{Delete}\((H,x)\) & remove \(x\) & \(O(n)\) & \(O(\log n)\)\\[2pt]
\multicolumn{4}{@{}l}{\emph{Auxiliary}}\\
\textsc{Link}\((H,y,x)\) & make root \(y\) a child of root \(x\) & \(O(1)\) & --\\
\textsc{Consolidate}\((H)\) & link roots until all root degrees differ & \(O(\Dn+t(H))\) & --\\
\textsc{Cut}\((H,x,y)\) & move child \(x\) of \(y\) to the root list & \(O(1)\) & --\\
\textsc{Cascading-Cut}\((H,y)\) & cut marked ancestors, mark the first unmarked one & \(O(\text{height})\) & --\\
\textsc{Clear}\((H)\) & empty the heap, keep the node pool & \(O(1)\) & \(O(1)\)\\\bottomrule
\end{tabular}
\caption{The ADT implemented in \texttt{include/fib\_heap.hpp}. \(\Dn=\lfloor\log_\varphi n\rfloor\) bounds the
degree of any node (\cref{cor:degree}).}\label{tab:adt}
\end{table}

\subsection{Pseudocode}
The pseudocode follows the implementation line by line. One detail differs from the textbook version: when
\textsc{Extract-Min} moves the children of the minimum to the root list it also clears their marks. A root never
needs a mark (marks only matter for nodes that have a parent), and clearing them lowers the potential, so every bound
below still holds.

\begin{algorithm}[H]
\Fn{\textsc{Insert}\((H,x)\)}{
  \(x.\mathit{degree}\gets0;\ x.\mathit{parent}\gets x.\mathit{child}\gets\textsc{nil};\ x.\mathit{mark}\gets\textsc{false}\)\;
  add \(x\) to the root list of \(H\)\;
  \lIf{\(H.\mathit{min}=\textsc{nil}\) \textbf{or} \(x.\mathit{key}<H.\mathit{min}.\mathit{key}\)}{\(H.\mathit{min}\gets x\)}
  \(H.n\gets H.n+1\)\;
}
\Fn{\textsc{Union}\((H_1,H_2)\)}{
  splice the root list of \(H_2\) into that of \(H_1\) \tcp*{four pointer updates}
  \lIf{\(H_2.\mathit{min}.\mathit{key}<H_1.\mathit{min}.\mathit{key}\)}{\(H_1.\mathit{min}\gets H_2.\mathit{min}\)}
  \(H_1.n\gets H_1.n+H_2.n\); \(H_2\) becomes empty\;
}
\caption{\textsc{Insert} and \textsc{Union}: lazy, no restructuring.}
\end{algorithm}

\begin{algorithm}[H]
\Fn{\textsc{Extract-Min}\((H)\)}{
  \(z\gets H.\mathit{min}\)\;
  \ForEach{child \(x\) of \(z\)}{ add \(x\) to the root list; \(x.\mathit{parent}\gets\textsc{nil}\); \(x.\mathit{mark}\gets\textsc{false}\)\; }
  remove \(z\) from the root list; \(H.n\gets H.n-1\)\;
  \lIf{\(H.n=0\)}{\(H.\mathit{min}\gets\textsc{nil}\)}
  \lElse{\(H.\mathit{min}\gets z.\mathit{right}\); \textsc{Consolidate}\((H)\)}
  \Return \(z\)\;
}
\Fn{\textsc{Consolidate}\((H)\)}{
  \(A[0..\Dn+1]\gets\textsc{nil}\)\;
  \ForEach{root \(w\) in a snapshot of the root list}{
    \(x\gets w\); \(d\gets x.\mathit{degree}\)\;
    \While{\(A[d]\neq\textsc{nil}\)}{
      \(y\gets A[d]\); \lIf{\(y.\mathit{key}<x.\mathit{key}\)}{swap \(x,y\)}
      \textsc{Link}\((H,y,x)\); \(A[d]\gets\textsc{nil}\); \(d\gets d+1\)\;
    }
    \(A[d]\gets x\)\;
  }
  rebuild the root list from the non-empty entries of \(A\) and set \(H.\mathit{min}\)\;
}
\Fn{\textsc{Link}\((H,y,x)\)}{
  remove \(y\) from the root list; make \(y\) a child of \(x\); \(x.\mathit{degree}\gets x.\mathit{degree}+1\); \(y.\mathit{mark}\gets\textsc{false}\)\;
}
\caption{\textsc{Extract-Min}: all deferred work happens here.}
\end{algorithm}

\begin{algorithm}[H]
\Fn{\textsc{Decrease-Key}\((H,x,k)\)}{
  \lIf{\(k>x.\mathit{key}\)}{\textbf{error}}
  \(x.\mathit{key}\gets k\); \(y\gets x.\mathit{parent}\)\;
  \If{\(y\neq\textsc{nil}\) \textbf{and} \(x.\mathit{key}<y.\mathit{key}\)}{ \textsc{Cut}\((H,x,y)\); \textsc{Cascading-Cut}\((H,y)\)\; }
  \lIf{\(x.\mathit{key}<H.\mathit{min}.\mathit{key}\)}{\(H.\mathit{min}\gets x\)}
}
\Fn{\textsc{Cut}\((H,x,y)\)}{
  remove \(x\) from the child list of \(y\); \(y.\mathit{degree}\gets y.\mathit{degree}-1\)\;
  add \(x\) to the root list; \(x.\mathit{parent}\gets\textsc{nil}\); \(x.\mathit{mark}\gets\textsc{false}\)\;
}
\Fn{\textsc{Cascading-Cut}\((H,y)\)}{
  \While{\(z\gets y.\mathit{parent}\neq\textsc{nil}\)}{
    \lIf{\(y.\mathit{mark}=\textsc{false}\)}{\(y.\mathit{mark}\gets\textsc{true}\); \textbf{return}}
    \textsc{Cut}\((H,y,z)\); \(y\gets z\)\;
  }
}
\Fn{\textsc{Delete}\((H,x)\)}{ \textsc{Decrease-Key}\((H,x,-\infty)\); \textsc{Extract-Min}\((H)\)\; }
\caption{\textsc{Decrease-Key}, the cut rules and \textsc{Delete}. \textsc{Cascading-Cut} is written as a loop, as in the code.}
\end{algorithm}

\subsection{A worked example}
\Cref{fig:heapex} runs \textsc{Extract-Min} and one \textsc{Decrease-Key} on a small heap. Before the extraction
there are five roots; afterwards three, with degrees 3, 1 and 0. The decrease-key on 35 then triggers a cascading cut
because its parent 26 is already marked.

\begin{figure}[H]\centering
\tikzset{fn/.style={circle,draw,minimum size=6.2mm,inner sep=0pt,font=\scriptsize},
         mk/.style={fn,fill=black!70,text=white}, mn/.style={fn,very thick,draw=accent},
         level distance=9mm, sibling distance=8mm, edge from parent/.style={draw}}
\begin{subfigure}{\linewidth}\centering
\begin{tikzpicture}
\node[fn] at (0,0) {7};
\node[mn] at (1.1,0) {3} child{node[fn]{18} child{node[fn]{39}}} child{node[fn]{52}};
\node[fn] at (2.6,0) {17} child{node[fn]{30}};
\node[fn] at (3.9,0) {24} child{node[mk]{26} child{node[fn]{35}}} child{node[fn]{46}};
\node[fn] at (5.2,0) {23};
\draw[densely dotted] (-0.4,0.45) -- (5.6,0.45) node[right,font=\scriptsize]{root list};
\end{tikzpicture}
\caption{Before: \(t(H)=5\), \(m(H)=1\) (26 is marked), \(\Phi=7\). The minimum 3 is outlined.}
\end{subfigure}\\[4pt]
\begin{subfigure}{\linewidth}\centering
\begin{tikzpicture}[sibling distance=12mm]
\node[mn] at (0,0) {7} child{node[fn]{23}} child{node[fn]{17} child{node[fn]{30}}}
  child[sibling distance=14mm]{node[fn]{24} child[sibling distance=8mm]{node[mk]{26} child{node[fn]{35}}} child[sibling distance=8mm]{node[fn]{46}}};
\node[fn] at (2.7,0) {18} child{node[fn]{39}};
\node[fn] at (3.8,0) {52};
\end{tikzpicture}
\caption{After \textsc{Extract-Min}: 18 and 52 are promoted; \textsc{Consolidate} links 23, 17 and 24 under 7
(degrees 0, 1, 2 become one tree of degree 3). \(t(H)=3\), \(\Phi=5\).}
\end{subfigure}\\[4pt]
\begin{subfigure}{\linewidth}\centering
\begin{tikzpicture}[sibling distance=12mm]
\node[fn] at (0,0) {7} child{node[fn]{23}} child{node[fn]{17} child{node[fn]{30}}}
  child{node[mk]{24} child{node[fn]{46}}};
\node[fn] at (2.2,0) {18} child{node[fn]{39}};
\node[fn] at (3.2,0) {52};
\node[mn] at (4.2,0) {5};
\node[fn] at (5.2,0) {26};
\end{tikzpicture}
\caption{After \textsc{Decrease-Key}(35, 5): 35 is cut; its parent 26 was marked, so 26 is cut too (cascading cut);
24 had not lost a child before, so it is marked and the cascade stops. Two cuts, \(\Phi\) goes from 5 to 7.}
\end{subfigure}
\caption{\textsc{Extract-Min} and a cascading cut on a small Fibonacci heap. Filled nodes are marked.}\label{fig:heapex}
\end{figure}

% ------------------------------------------------------------------ 3
\section{Amortized analysis}\label{sec:amortized}

A single \textsc{Extract-Min} can take \(\Theta(n)\) time (after \(n\) inserts the root list has \(n\) trees), and a
single \textsc{Decrease-Key} can cut a whole path. Worst-case bounds per operation are therefore useless for
Dijkstra. What matters is the total cost of a \emph{sequence} of operations, and that is what amortized analysis
bounds. We use the potential method and then restate the result with the accounting and aggregate methods.

\subsection{Cost model and potential function}
Let \(D_0\) be the empty heap and \(D_i\) the heap after the \(i\)-th operation, whose actual cost is \(c_i\). For a
potential \(\Phi\) with \(\Phi(D_0)=0\) and \(\Phi(D_i)\ge0\), the amortized cost is
\[\hat c_i = c_i+\Phi(D_i)-\Phi(D_{i-1}),\qquad\text{so}\qquad
\sum_{i=1}^{N}\hat c_i=\sum_{i=1}^{N}c_i+\Phi(D_N)-\Phi(D_0)\;\ge\;\sum_{i=1}^{N}c_i .\]
The sum of amortized costs is an upper bound on the true total for every sequence.

\begin{definition}[potential]\label{def:phi}
\(\Phi(H)=t(H)+2\,m(H)\), where \(t(H)\) is the number of trees in the root list and \(m(H)\) the number of marked nodes.
\end{definition}

\noindent\(\Phi(D_0)=0\) and \(\Phi\ge0\) always. One unit of cost pays for a constant amount of pointer work. We
charge \textsc{Insert} and \textsc{Union} one unit each; \textsc{Decrease-Key} one unit plus one per cut; and
\textsc{Extract-Min} one unit, plus one per promoted child, plus one per root examined by \textsc{Consolidate}.
Each \textsc{Link} and each iteration of the inner \textbf{while} loop removes one root that was examined, so
there are at most as many of them as examined roots, and the true running time is within a constant factor of this
count. The implementation counts exactly these units (\texttt{FibStats::actual\_cost}).

\subsection{Amortized cost of each operation}

\begin{lemma}[\textsc{Insert}, \textsc{Union}, \textsc{Make-Heap}, \textsc{Minimum}]\label{lem:insert}
\(\hat c_{\textsc{Insert}}=2\), \(\hat c_{\textsc{Union}}=1\), and \(\hat c=1\) for \textsc{Make-Heap} and \textsc{Minimum}.
\end{lemma}
\begin{proof}
\textsc{Insert} adds one unmarked root: \(c=1\), \(\Delta t=1\), \(\Delta m=0\), so \(\hat c=1+1=2\).
\textsc{Union} concatenates two root lists: \(t(H_1\cup H_2)=t(H_1)+t(H_2)\) and the same holds for \(m\), so
\(\Phi\) of the result equals \(\Phi(H_1)+\Phi(H_2)\), \(\Delta\Phi=0\) and \(\hat c=1\).
The other two change nothing.
\end{proof}

\begin{lemma}[\textsc{Extract-Min}]\label{lem:extract}
If the heap has \(n\) nodes and every node has degree at most \(\Dn\), then \(\hat c_{\textsc{Extract-Min}}\le 3\Dn+1\).
\end{lemma}
\begin{proof}
Let \(t=t(H)\) before the call and let \(c'\le\Dn\) be the number of children of the minimum \(z\). Promoting the
children costs \(c'\). After \(z\) is removed the root list has \(t+c'-1\) trees and \textsc{Consolidate} examines each,
so
\[c_i=1+c'+(t+c'-1)=2c'+t .\]
After \textsc{Consolidate} no two roots have the same degree and every degree is in \(\{0,\dots,\Dn\}\), so at most
\(\Dn+1\) roots remain. No node becomes marked (\textsc{Link} and the promotion only clear marks), so
\(\Delta m\le0\) and
\[\hat c_i\le 2c'+t+\big((\Dn+1)-t\big)=2c'+\Dn+1\le3\Dn+1. \qedhere\]
\end{proof}
The term \(t\) cancels: the long root list built up by lazy inserts was paid for in advance by the potential those
inserts left behind.

\begin{lemma}[\textsc{Decrease-Key}]\label{lem:decrease}
\(\hat c_{\textsc{Decrease-Key}}\le5\).
\end{lemma}
\begin{proof}
If no cut happens, \(c_i=1\) and \(\Delta\Phi=0\). Otherwise let \(c\ge1\) be the number of cuts: the cut of \(x\)
and \(c-1\) cascading cuts. Then \(c_i=1+c\). Each cut adds one root, \(\Delta t=c\). Each cascading cut cuts a marked
node and clears its mark, the cut of \(x\) clears its mark if it had one, and at most one node (where the cascade
stops) becomes marked, so \(\Delta m\le 1-(c-1)=2-c\). Hence
\[\hat c_i\le(1+c)+c+2(2-c)=5. \qedhere\]
\end{proof}
This is why marked nodes carry weight two in \(\Phi\): one unit pays for the cut of the node and one for the
new root it creates.

\begin{lemma}[\textsc{Delete}]\(\hat c_{\textsc{Delete}}\le 3\Dn+6\), since it is one \textsc{Decrease-Key} followed by one \textsc{Extract-Min}.\end{lemma}

\subsection{Bounding the maximum degree}
The bound \(\Dn=O(\log n)\) is what makes \cref{lem:extract} logarithmic, and it depends on the marking rule.

\begin{lemma}\label{lem:children}
Let \(x\) be a node and \(y_1,\dots,y_d\) its children in the order in which they were linked to \(x\). Then
\(y_i.\mathit{degree}\ge i-2\) for \(i=1,\dots,d\).
\end{lemma}
\begin{proof}
When \(y_i\) was linked, \(y_1,\dots,y_{i-1}\) were already children of \(x\), so \(x\) had degree at least \(i-1\).
\textsc{Link} joins roots of equal degree, so \(y_i\) also had degree at least \(i-1\). Since then \(y_i\) has lost at
most one child: the first loss marks it, and a second loss would have cut \(y_i\) from \(x\). So its degree is at
least \(i-2\).
\end{proof}

\begin{lemma}\label{lem:size}
A node of degree \(d\) roots a subtree with at least \(F_{d+2}\ge\varphi^{d}\) nodes, where \(F_j\) is the
\(j\)-th Fibonacci number and \(\varphi=(1+\sqrt5)/2\).
\end{lemma}
\begin{proof}
Let \(s_d\) be the smallest possible subtree size of a node of degree \(d\); \(s_0=1\), \(s_1=2\), and \(s_d\) is
non-decreasing in \(d\). By \cref{lem:children}, counting \(x\), its first child and the subtrees of the others,
\(s_d\ge 2+\sum_{i=2}^{d}s_{i-2}\). Induction with \(F_{d+2}=1+\sum_{j=0}^{d}F_j\) gives \(s_d\ge F_{d+2}\), and
\(F_{d+2}\ge\varphi^{d}\) follows from \(\varphi^2=\varphi+1\) by induction.
\end{proof}

\begin{corollary}\label{cor:degree}
Every node of an \(n\)-node Fibonacci heap has degree at most \(\Dn=\lfloor\log_\varphi n\rfloor\le1.4405\log_2 n\).
\end{corollary}
\noindent The array \(A\) in \textsc{Consolidate} is sized by this bound, and the code asserts it on every link.

\subsection{The three methods side by side}
\begin{theorem}\label{thm:sequence}
Starting from an empty heap, any sequence of \(a\) \textsc{Insert}/\textsc{Union} operations, \(b\) \textsc{Extract-Min}
operations, \(c\) \textsc{Decrease-Key} operations and \(d\) \textsc{Delete} operations, in which the heap never holds more
than \(n\) nodes, has total actual cost at most \(2a+b\,(3\Dn+1)+5c+d\,(3\Dn+6)=O(a+c+(b+d)\log n)\).
\end{theorem}
\begin{proof}
Sum the bounds of \cref{lem:insert,lem:extract,lem:decrease} and use \(\sum c_i\le\sum\hat c_i\).
\end{proof}

\paragraph{Accounting method.} The same analysis can be told with credits. Every root holds one credit and every
marked node two. \textsc{Insert} is charged 2: one pays for the work, one is left on the new root. \textsc{Consolidate}
pays for examining and linking a root with that root's credit. \textsc{Decrease-Key} is charged 5: a cascading cut is
paid by the two credits of the marked node it cuts (one for the cut, one placed on the new root), and the charge covers
the first cut, its root credit, and the two credits for the one new mark. The credit balance is exactly
\(\Phi(H)\), so it never goes negative.

\paragraph{Aggregate method for Dijkstra and Prim.} A Dijkstra run on \(n\) vertices and \(m\) edges performs at most
\(n\) \textsc{Insert}, \(n\) \textsc{Extract-Min} and \(m\) \textsc{Decrease-Key} (an edge can lower a label only when its
first endpoint is extracted). By \cref{thm:sequence} the total heap work is at most
\(2n+n(3\Dn+1)+5m=O(m+n\log n)\), the Fredman--Tarjan bound. The heap is emptied at the end, so the next run starts
from \(\Phi=0\) and the bounds add up over the runs of a solver. For Prim on the complete distance network \(D\)
(G1, G3) the counts are \(k\), \(k\) and \(\le k^2\), giving \(O(k^2)\) instead of \(O(k^2\log k)\) with a binary heap.

\subsection{Checking the analysis on real runs}\label{sec:amortized-check}
The proofs above are checked in two ways.

\textbf{Operation by operation.} \texttt{tests/test\_fib\_heap.cpp} runs five random sequences of \(2\times10^5\)
operations (40\% \textsc{Insert}, 30\% \textsc{Decrease-Key}, 25\% \textsc{Extract-Min}, 5\% \textsc{Delete}), compares every result with
\texttt{std::multiset}, verifies the heap invariants every 4096 operations, and asserts after every single operation
that \(\Phi\ge0\) and \(c_i+\Delta\Phi\) does not exceed the bound of its lemma. All \(10^6\) checks pass.
\Cref{fig:amort} plots the running totals: the true cost stays at about \val{amort.test.ratio} of the amortized bound.

\textbf{Inside the solvers.} Each solver reports the heap counters of its run. \Cref{tab:heapcounts} sums them over
all runs of the main benchmark and compares the measured work with the bound of \cref{thm:sequence}.

\begin{figure}[!htbp]\centering
\begin{tikzpicture}
\begin{axis}[width=0.62\linewidth,height=5.2cm,xlabel={operations},ylabel={cumulative cost (units)},
  legend pos=north west,legend style={font=\footnotesize,draw=none},tick label style={font=\footnotesize},
  label style={font=\small},scaled ticks=true,grid=major,grid style={black!8}]
\addplot[cBB,thick,dashed] table[x=op,y=sum_amortized_bound,col sep=comma]{generated/amort_seed1.csv};
\addlegendentry{\(\sum\hat c_i\) (bound)}
\addplot[cG3,thick] table[x=op,y=sum_actual,col sep=comma]{generated/amort_seed1.csv};
\addlegendentry{\(\sum c_i\) (measured)}
\addplot[cG1,thin] table[x=op,y=potential,col sep=comma]{generated/amort_seed1.csv};
\addlegendentry{\(\Phi(D_i)\)}
\end{axis}
\end{tikzpicture}
\caption{Cumulative measured cost against the cumulative amortized bound on one random sequence of
\(2\times10^5\) operations (seed 1). The gap is the slack in the constants of \cref{lem:insert,lem:extract,lem:decrease};
\(\Phi\) stays small because \textsc{Extract-Min} keeps the root list short.}\label{fig:amort}
\end{figure}

\begin{table}[H]\centering\small
\begin{tabular}{@{}lrrrrrrrrr@{}}\toprule
Solver & runs & \textsc{Insert} & \textsc{Extr.-Min} & \textsc{Decr.-Key} & links & cuts & actual & bound & ratio\\\midrule
G1 & 199 & 1.2M & 1.2M & 0.6M & 5.0M & 0.16M & 17.8M & 50.2M & 0.35\\
G2 & 199 & 139.9M & 74.2M & 68.6M & 417.2M & 22.51M & 1380.2M & 3345.5M & 0.41\\
G3 & 199 & 1.3M & 1.3M & 0.9M & 5.3M & 0.26M & 18.9M & 53.9M & 0.35\\
DP & 174 & 1643.5M & 1643.5M & 788.0M & 5957.1M & 165.95M & 21044.4M & 53813.1M & 0.39\\
B\&B & 199 & 1708.3M & 868.2M & 493.8M & 4450.4M & 132.18M & 14744.0M & 60317.1M & 0.24\\
\bottomrule\end{tabular}
\caption{Heap work summed over all runs of the main benchmark (M = million). \emph{actual} is the measured cost in the units of \cref{sec:amortized}; \emph{bound} is \(2\,\#\textsc{Insert}+(3D(N)+1)\,\#\textsc{Extract-Min}+5\,\#\textsc{Decrease-Key}\) from \cref{thm:sequence}, with \(N\) the largest possible heap of the run. For B\&B the counts cover the vertex queue and the open-node queue. On these sparse benchmark graphs there are fewer decrease-keys than extract-mins in every solver; \cref{sec:ablation} shows what this means for running time.}\label{tab:heapcounts}
\end{table}


% ------------------------------------------------------------------ 4
\section{Three greedy strategies}\label{sec:greedy}

The three strategies differ in the greedy \emph{selection criterion}, that is, in which object they commit to at each step:
\begin{description}[leftmargin=1.2cm,style=nextline]
\item[G1] the cheapest shortest path between two terminal components (an MST of the distance network \(D\));
\item[G2] the cheapest shortest path from the single growing tree to a terminal outside it;
\item[G3] the triple of terminals whose connection through one Steiner vertex saves the most against the current
spanning tree.
\end{description}
G1 and G2 only ever join terminals in pairs. G3 is the only one that can choose a Steiner vertex because it serves
three terminals at once, and \cref{sec:gaps} shows that this is the step that pays off. All three end with the same
clean-up, which we prove safe first.

\begin{lemma}[clean-up]\label{lem:cleanup}
Let \(H\subseteq G\) be connected with \(T\subseteq V(H)\). Replacing \(H\) by a minimum spanning tree of
\(G[V(H)]\) and then deleting non-terminal leaves until none is left produces a Steiner tree of cost at most \(w(H)\).
\end{lemma}
\begin{proof}
\(H\) contains a spanning tree of \(V(H)\) of cost at most \(w(H)\) (weights are non-negative), and \(H\subseteq G[V(H)]\), so
the MST of \(G[V(H)]\) costs at most \(w(H)\). Deleting a non-terminal leaf keeps a tree that still contains \(T\) and removes
an edge of weight \(\ge0\). The clean-up runs Prim on the heap (\(O(m+n\log n)\)) and a leaf queue (\(O(n)\)).
\end{proof}

\subsection{G1: distance-network MST (Kou, Markowsky and Berman)}
\begin{algorithm}[H]
\For{\(i=1\) \KwTo \(k\)}{ Dijkstra from \(t_i\) (Fibonacci heap); store \(d(t_i,\cdot)\) and the predecessor tree\; }
build \(D\) with \(D_{ij}=d(t_i,t_j)\); \(M\gets\) Prim on \(D\) (Fibonacci heap)\;
\(H\gets\) union of the shortest paths of \(G\) for the edges of \(M\)\;
\Return \textsc{Clean-Up}\((H)\)\;
\caption{G1 (\texttt{src/greedy1\_kmb.cpp}).}
\end{algorithm}

\paragraph{Running time.} \(k\) Dijkstra runs cost \(O(k(m+n\log n))\) by \cref{sec:amortized}; Prim on \(D\) costs \(O(k^2)\);
expanding \(k-1\) paths costs \(O(kn)\); the clean-up costs \(O(m+n\log n)\). Total \(O(k(m+n\log n))\), polynomial.
The original paper states \(O(kn^2)\) because it uses an array-based Dijkstra.

\begin{theorem}\label{thm:kmb}
Let \(\ell\ge2\) be the number of leaves of an optimal Steiner tree \(S^\ast\). Then
\(w(\mathrm{MST}(D))\le\big(2-\tfrac{2}{\ell}\big)\OPT\), and G1 returns a tree of cost at most \(\big(2-\tfrac{2}{\ell}\big)\OPT\le\big(2-\tfrac2k\big)\OPT\).
\end{theorem}
\begin{proof}
Double every edge of \(S^\ast\). The result is Eulerian; an Euler tour \(W\) has cost \(2\,\OPT\) and visits the leaves
of \(S^\ast\) in some cyclic order \(L_1,\dots,L_\ell\). Cut \(W\) at the leaves into \(\ell\) walks \(W_j\) from \(L_j\) to
\(L_{j+1}\). The longest of them costs at least \(2\OPT/\ell\); delete it. The other \(\ell-1\) walks form one walk of cost at
most \(2\OPT(1-1/\ell)\) that still visits every terminal: a leaf is an end of two walks, and an inner vertex of degree
\(\delta\ge2\) appears in \(\delta\) walks, so deleting one walk removes neither. Listing the terminals in the order this walk
first reaches them gives a path in \(D\) whose edges are no longer than the walk segments between them (\(d\) is a shortest-path
metric). That path is a spanning tree of \(D\), so \(w(\mathrm{MST}(D))\le(2-2/\ell)\OPT\). The expanded subgraph \(H\) costs at most
\(w(\mathrm{MST}(D))\) (shared edges are counted once), and \cref{lem:cleanup} does not increase it. Since \(\ell\le k\), the bound
\(2-2/k\) follows.
\end{proof}

\subsection{G2: shortest-path heuristic (Takahashi and Matsuyama)}
\begin{algorithm}[H]
\(S\gets\{r\}\)\;
\While{\(T\not\subseteq V(S)\)}{
  multi-source Dijkstra with label 0 on every vertex of \(S\), stopped when the first terminal \(t\notin V(S)\) is extracted\;
  add the shortest path from \(S\) to \(t\) to \(S\)\;
}
\Return \textsc{Clean-Up}\((S)\)\;
\caption{G2 for one root \(r\) (\texttt{src/greedy2\_sph.cpp}). The solver runs every terminal as the root and keeps the best tree.}
\end{algorithm}

\paragraph{Running time.} Each of the \(k-1\) rounds is one Dijkstra, \(O(m+n\log n)\), and adding a path is \(O(n)\). One root
costs \(O(k(m+n\log n))\); all \(k\) roots cost \(O(k^2(m+n\log n))\). Stopping at the first terminal is safe because a Dijkstra
label is final once it is extracted, and it shortens the rounds in practice.

\begin{theorem}\label{thm:sph}
For every root, G2 returns a tree of cost at most \(w(\mathrm{MST}(D))\le\big(2-\tfrac2\ell\big)\OPT\).
\end{theorem}
\begin{proof}
Let \(S_0=\{r\}\subset S_1\subset\dots\subset S_{k-1}=T\) be the terminal sets after each round and \(\delta_i\) the cost of the
\(i\)-th path. Because the path goes from the current tree, which contains \(S_{i-1}\), to the nearest outside terminal,
\(\delta_i\le\min\{d(s,t):s\in S_{i-1},\,t\notin S_{i-1}\}\), the cheapest edge of \(D\) crossing the cut \((S_{i-1},T\setminus S_{i-1})\).
Let \(M\) be an MST of \(D\). We match round \(i\) to a distinct edge of \(M\) that crosses cut \(i\). For any \(j\) rounds
\(i_1<\dots<i_j\) the cuts split \(T\) into \(j+1\) non-empty layers \(S_{i_1-1},\,S_{i_2-1}\setminus S_{i_1-1},\dots,T\setminus S_{i_j-1}\).
\(M\) is connected, so at least \(j\) of its edges run between different layers, and each of them crosses at least one of the
\(j\) cuts. This is Hall's condition, so the matching exists and \(\sum_i\delta_i\le w(M)\). The clean-up does not increase the
cost, and \cref{thm:kmb} bounds \(w(M)\).
\end{proof}

\subsection{G3: best-gain triples (Zelikovsky)}
For a triple \(z=\{a,b,c\}\subseteq T\) let \(v(z)=\arg\min_v d(a,v)+d(b,v)+d(c,v)\) be its best centre and \(d(z)\) that
minimum (the cheapest star joining \(a,b,c\)). Let \(F\) be the current distance network and \(F[z]\) the network in which
\(a,b\) and \(b,c\) are joined by zero-cost edges. The \emph{win} of \(z\) is
\[\mathrm{win}_F(z)=w(\mathrm{MST}(F))-w(\mathrm{MST}(F[z]))-d(z).\]

\begin{algorithm}[H]
Dijkstra from every terminal; \(F\gets D\); \(Z\gets\emptyset\)\;
\Repeat{\(\max_z\mathrm{win}_F(z)\le0\)}{
  \(M\gets\) Prim on \(F\); compute \(s_{xy}\) = heaviest edge on the \(M\)-path between \(x,y\), for all pairs\;
  \(z^\ast\gets\arg\max_z\mathrm{win}_F(z)\), using \cref{lem:gain} and lazy evaluation of \(d(z)\)\;
  \lIf{\(\mathrm{win}_F(z^\ast)>0\)}{\(F\gets F[z^\ast]\); \(Z\gets Z\cup\{z^\ast\}\)}
}
\(H\gets\) stars of all \(z\in Z\) \(\cup\) shortest paths of the non-zero edges of \(\mathrm{MST}(F)\)\;
\Return \textsc{Clean-Up}\((H)\)\;
\caption{G3 (\texttt{src/greedy3\_zelikovsky.cpp}).}
\end{algorithm}

\begin{lemma}\label{lem:gain}
Let \(M\) be an MST of \(F\), and let \(s_{ab},s_{bc},s_{ac}\) be the heaviest edge weights on the \(M\)-paths between the three
pairs. Then \(w(\mathrm{MST}(F))-w(\mathrm{MST}(F[z]))=\max(s_{ab},s_{bc},s_{ac})+\min(s_{ab},s_{bc},s_{ac})\).
\end{lemma}
\begin{proof}
The three \(M\)-paths form a subtree shaped like a Y with centre \(o\) (one leg may be empty). Let \(\mu_a,\mu_b,\mu_c\) be the
heaviest weights on the legs from \(o\) to \(a,b,c\) (0 for an empty leg). Joining \(a,b,c\) by zero edges is the same as
contracting them into one node, which turns the Y into two independent cycles through \(o\). An MST of \(F[z]\) is obtained from
\(M\) plus the zero edges by deleting one edge on each cycle, and the two deleted edges must lie on two different legs (deleting
two edges of one leg disconnects the part between them). The best choice deletes the heaviest edge of each of the two legs with the
largest maxima, saving the sum of the two largest of \(\mu_a,\mu_b,\mu_c\). Since \(s_{ab}=\max(\mu_a,\mu_b)\) and similarly for the
other pairs, if \(\mu_a\ge\mu_b\ge\mu_c\) the three values are \(\mu_a,\mu_b,\mu_a\), and max plus min is \(\mu_a+\mu_b\). The other
orders are symmetric.
\end{proof}

\paragraph{Running time.} \(k\) Dijkstra runs cost \(O(k(m+n\log n))\). Each round costs \(O(k^2)\) for Prim on the heap and for the
\(s\) matrix (one traversal of \(M\) per terminal), \(O(k^3)\) to scan the triples with \cref{lem:gain}, and \(O(n)\) per centre that
has to be computed. Centres are cached across rounds, so each of the \(\binom{k}{3}\) triples has its centre computed at most once,
\(O(k^3n)\) in all. Every accepted triple merges three components into one, so there are at most \((k-1)/2\) rounds, and the total
is \(O(k(m+n\log n)+k^4+k^3n)\). In practice few centres are computed: a triple is skipped when its MST saving alone cannot beat the
best win so far, or when it cannot beat it even with the lower bound \(d(z)\ge(D_{ab}+D_{bc}+D_{ac})/2\) (each leg of the star lies on two of
the three pairwise paths). On \texttt{b18} (\(k=50\)) the code scans \val{g3.b18.scanned} triples and computes \val{g3.b18.centres} centres.

\begin{theorem}[Zelikovsky 1993]\label{thm:zel}
G3 returns a tree of cost at most \(w(\mathrm{MST}(D))\le(2-2/\ell)\OPT\), and at most \(\tfrac{11}{6}\OPT\).
\end{theorem}
\begin{proof}[Proof of the first bound and outline of the second]
Write \(C(F,Z)=w(\mathrm{MST}(F))+\sum_{z\in Z}d(z)\). At the start \(C=w(\mathrm{MST}(D))\), and each accepted triple lowers \(C\)
by its win, which is positive. The final \(H\) costs at most \(C\) and \cref{lem:cleanup} keeps it there, so the output costs at most
\(w(\mathrm{MST}(D))\). For \(11/6\), Zelikovsky decomposes an optimal tree into full components and shows that a best 3-restricted
Steiner tree, \(\mathrm{smt}_3\), costs at most \(\tfrac53\OPT\). A charging argument over the triples of that tree then shows that
the greedy choice of the best-win triple gives
\(\mathrm{cost}\le\tfrac12\big(w(\mathrm{MST}(D))+\mathrm{smt}_3\big)\le\tfrac12\big(2+\tfrac53\big)\OPT=\tfrac{11}{6}\OPT\).
The full proof is in \cite{zel93}; we do not reproduce it.
\end{proof}
\noindent The best known polynomial ratio is \(\ln4+\varepsilon\approx1.39\) \cite{bgrs}; it needs an LP solver, which the
assignment rules out.

\subsection{Counterexamples}\label{sec:counter}
\paragraph{A tight family for G1 and G2.}
Take a Steiner vertex \(c\) joined to \(\ell\) terminals by edges of weight 10, and join every pair of terminals by an edge of
weight 19 (\cref{fig:star}). The star costs \(10\ell\) and is optimal. Every terminal pair is at distance \(19<20\), so \(D\)
is the uniform 19-clique; G1 and G2 both build a path of \(\ell-1\) direct edges, and the clean-up keeps it because \(c\) is never
reached. Their cost is \(19(\ell-1)\), a ratio of \(1.9(1-1/\ell)\), which approaches the bound \(2-2/\ell\) of
\cref{thm:kmb,thm:sph} as \(19/10\) is replaced by \((2w-1)/w\) with larger \(w\). G3 is optimal on the whole family: every triple has
win \(2\cdot19-30=8>0\), so G3 contracts triples through \(c\), \(c\) enters \(H\), and the clean-up's MST then finds the star.
\Cref{tab:star} gives the measured costs; they match the formulas exactly.

\begin{figure}[H]\centering
\begin{minipage}{0.36\linewidth}\centering
\begin{tikzpicture}[scale=1.05]
\foreach \i in {1,...,6}{\node[circle,draw,fill=accent,text=white,inner sep=1.5pt,font=\scriptsize] (t\i) at ({60*\i}:1.6) {\(t_\i\)};}
\node[circle,draw,inner sep=1.5pt,font=\scriptsize] (c) at (0,0) {\(c\)};
\foreach \i in {1,...,6}{\draw[thick,cG3] (c) -- node[font=\tiny,fill=white,inner sep=0.5pt]{10} (t\i);}
\foreach \i/\j in {1/2,2/3,3/4,4/5,5/6,6/1,1/3,2/4,3/5,4/6,5/1,6/2,1/4,2/5,3/6}{\draw[black!25] (t\i) -- (t\j);}
\foreach \i/\j in {1/2,2/3,3/4,4/5,5/6}{\draw[very thick,cG1,dashed] (t\i) -- (t\j);}
\end{tikzpicture}
\end{minipage}\hfill
\begin{minipage}{0.6\linewidth}
\small\begin{tabular}{@{}rrrrrrr@{}}\toprule
\(\ell\) & OPT & G1 & G2 & G3 & G1/OPT & \(2-2/\ell\)\\\midrule
3 & 30 & 38 & 38 & 30 & 1.267 & 1.333\\
4 & 40 & 57 & 57 & 40 & 1.425 & 1.500\\
5 & 50 & 76 & 76 & 50 & 1.520 & 1.600\\
6 & 60 & 95 & 95 & 60 & 1.583 & 1.667\\
7 & 70 & 114 & 114 & 70 & 1.629 & 1.714\\
8 & 80 & 133 & 133 & 80 & 1.663 & 1.750\\
9 & 90 & 152 & 152 & 90 & 1.689 & 1.778\\
10 & 100 & 171 & 171 & 100 & 1.710 & 1.800\\
11 & 110 & 190 & 190 & 110 & 1.727 & 1.818\\
12 & 120 & 209 & 209 & 120 & 1.742 & 1.833\\
\bottomrule\end{tabular}
\captionof{table}{Measured costs on the star family.}\label{tab:star}

\end{minipage}
\caption{Left: the gadget for \(\ell=6\). Grey edges have weight 19. The dashed path is what G1 and G2 return (cost 95); the star
through \(c\) is optimal (cost 60) and is what G3 returns. Right: measured costs for \(\ell=3,\dots,12\).}\label{fig:star}
\end{figure}

\paragraph{Instances found by search.}
For G3 the star family is no counterexample, and the theorem gives no tight family that is easy to write down. So
\texttt{scripts/gen\_instances.py} searches for one. It starts from small random graphs (7 to 9 vertices, 4 or 5 terminals) and
hill-climbs on edge weights, edges and the terminal set to maximise \(\mathrm{cost}(\mathrm{G}x)/\OPT\), with \(\OPT\) from the
dynamic program, accepting ties so it can cross plateaus. It then deletes edges as long as the ratio does not drop, so that the
graphs stay small enough to draw. Six restarts per strategy, all seeded. \Cref{fig:gadgets} shows the worst instance found
for each strategy. The search found nothing above 1.4 on graphs this small, far from the proven bounds: those bounds need many leaves,
and a 4- or 5-terminal tree has few.

\begin{figure}[!htbp]\centering
\begin{subfigure}{0.32\linewidth}\centering\begin{tikzpicture}[scale=1.25]
\node[circle,draw,fill=white,inner sep=1pt,minimum size=4.6mm,font=\scriptsize] (v0) at (90.0:1.35) {1};
\node[circle,draw,fill=accent,text=white,inner sep=1pt,minimum size=4.6mm,font=\scriptsize] (v1) at (141.4:1.35) {2};
\node[circle,draw,fill=white,inner sep=1pt,minimum size=4.6mm,font=\scriptsize] (v2) at (192.9:1.35) {3};
\node[circle,draw,fill=accent,text=white,inner sep=1pt,minimum size=4.6mm,font=\scriptsize] (v3) at (244.3:1.35) {4};
\node[circle,draw,fill=accent,text=white,inner sep=1pt,minimum size=4.6mm,font=\scriptsize] (v4) at (295.7:1.35) {5};
\node[circle,draw,fill=accent,text=white,inner sep=1pt,minimum size=4.6mm,font=\scriptsize] (v5) at (347.1:1.35) {6};
\node[circle,draw,fill=white,inner sep=1pt,minimum size=4.6mm,font=\scriptsize] (v6) at (398.6:1.35) {7};
\draw[black!30] (v0) -- node[font=\tiny,fill=white,inner sep=0.4pt]{12} (v1);
\draw[cG1,thick,dashed] (v0) -- (v1);
\draw[line width=2.6pt,cG3!45] (v0) -- node[font=\tiny,fill=white,inner sep=0.4pt]{4} (v2);
\draw[line width=2.6pt,cG3!45] (v0) -- node[font=\tiny,fill=white,inner sep=0.4pt]{3} (v3);
\draw[cG1,thick,dashed] (v0) -- (v3);
\draw[line width=2.6pt,cG3!45] (v2) -- node[font=\tiny,fill=white,inner sep=0.4pt]{10} (v4);
\draw[line width=2.6pt,cG3!45] (v1) -- node[font=\tiny,fill=white,inner sep=0.4pt]{5} (v6);
\draw[cG1,thick,dashed] (v1) -- (v6);
\draw[black!30] (v4) -- node[font=\tiny,fill=white,inner sep=0.4pt]{17} (v5);
\draw[cG1,thick,dashed] (v4) -- (v5);
\draw[line width=2.6pt,cG3!45] (v2) -- node[font=\tiny,fill=white,inner sep=0.4pt]{3} (v6);
\draw[line width=2.6pt,cG3!45] (v5) -- node[font=\tiny,fill=white,inner sep=0.4pt]{5} (v6);
\draw[cG1,thick,dashed] (v5) -- (v6);
\end{tikzpicture}
\caption{G1: cost 42, OPT 30, ratio 1.400.}\end{subfigure}\hfill
\begin{subfigure}{0.32\linewidth}\centering\begin{tikzpicture}[scale=1.25]
\node[circle,draw,fill=accent,text=white,inner sep=1pt,minimum size=4.6mm,font=\scriptsize] (v0) at (90.0:1.35) {1};
\node[circle,draw,fill=white,inner sep=1pt,minimum size=4.6mm,font=\scriptsize] (v1) at (135.0:1.35) {2};
\node[circle,draw,fill=accent,text=white,inner sep=1pt,minimum size=4.6mm,font=\scriptsize] (v2) at (180.0:1.35) {3};
\node[circle,draw,fill=white,inner sep=1pt,minimum size=4.6mm,font=\scriptsize] (v3) at (225.0:1.35) {4};
\node[circle,draw,fill=white,inner sep=1pt,minimum size=4.6mm,font=\scriptsize] (v4) at (270.0:1.35) {5};
\node[circle,draw,fill=accent,text=white,inner sep=1pt,minimum size=4.6mm,font=\scriptsize] (v5) at (315.0:1.35) {6};
\node[circle,draw,fill=accent,text=white,inner sep=1pt,minimum size=4.6mm,font=\scriptsize] (v6) at (360.0:1.35) {7};
\node[circle,draw,fill=accent,text=white,inner sep=1pt,minimum size=4.6mm,font=\scriptsize] (v7) at (405.0:1.35) {8};
\draw[line width=2.6pt,cG3!45] (v3) -- node[font=\tiny,fill=white,inner sep=0.4pt]{10} (v6);
\draw[black!30] (v4) -- node[font=\tiny,fill=white,inner sep=0.4pt]{18} (v6);
\draw[cG1,thick,dashed] (v4) -- (v6);
\draw[line width=2.6pt,cG3!45] (v6) -- node[font=\tiny,fill=white,inner sep=0.4pt]{1} (v7);
\draw[cG1,thick,dashed] (v6) -- (v7);
\draw[line width=2.6pt,cG3!45] (v2) -- node[font=\tiny,fill=white,inner sep=0.4pt]{18} (v3);
\draw[black!30] (v2) -- node[font=\tiny,fill=white,inner sep=0.4pt]{9} (v4);
\draw[cG1,thick,dashed] (v2) -- (v4);
\draw[line width=2.6pt,cG3!45] (v3) -- node[font=\tiny,fill=white,inner sep=0.4pt]{10} (v5);
\draw[line width=2.6pt,cG3!45] (v0) -- node[font=\tiny,fill=white,inner sep=0.4pt]{10} (v3);
\draw[black!30] (v5) -- node[font=\tiny,fill=white,inner sep=0.4pt]{20} (v6);
\draw[cG1,thick,dashed] (v5) -- (v6);
\draw[black!30] (v0) -- node[font=\tiny,fill=white,inner sep=0.4pt]{20} (v7);
\draw[cG1,thick,dashed] (v0) -- (v7);
\draw[black!30] (v1) -- node[font=\tiny,fill=white,inner sep=0.4pt]{6} (v7);
\end{tikzpicture}
\caption{G2: cost 68, OPT 49, ratio 1.388.}\end{subfigure}\hfill
\begin{subfigure}{0.32\linewidth}\centering\begin{tikzpicture}[scale=1.25]
\node[circle,draw,fill=accent,text=white,inner sep=1pt,minimum size=4.6mm,font=\scriptsize] (v0) at (90.0:1.35) {1};
\node[circle,draw,fill=white,inner sep=1pt,minimum size=4.6mm,font=\scriptsize] (v1) at (130.0:1.35) {2};
\node[circle,draw,fill=accent,text=white,inner sep=1pt,minimum size=4.6mm,font=\scriptsize] (v2) at (170.0:1.35) {3};
\node[circle,draw,fill=accent,text=white,inner sep=1pt,minimum size=4.6mm,font=\scriptsize] (v3) at (210.0:1.35) {4};
\node[circle,draw,fill=accent,text=white,inner sep=1pt,minimum size=4.6mm,font=\scriptsize] (v4) at (250.0:1.35) {5};
\node[circle,draw,fill=white,inner sep=1pt,minimum size=4.6mm,font=\scriptsize] (v5) at (290.0:1.35) {6};
\node[circle,draw,fill=white,inner sep=1pt,minimum size=4.6mm,font=\scriptsize] (v6) at (330.0:1.35) {7};
\node[circle,draw,fill=accent,text=white,inner sep=1pt,minimum size=4.6mm,font=\scriptsize] (v7) at (370.0:1.35) {8};
\node[circle,draw,fill=white,inner sep=1pt,minimum size=4.6mm,font=\scriptsize] (v8) at (410.0:1.35) {9};
\draw[line width=2.6pt,cG3!45] (v4) -- node[font=\tiny,fill=white,inner sep=0.4pt]{1} (v7);
\draw[cG1,thick,dashed] (v4) -- (v7);
\draw[black!30] (v1) -- node[font=\tiny,fill=white,inner sep=0.4pt]{11} (v7);
\draw[cG1,thick,dashed] (v1) -- (v7);
\draw[line width=2.6pt,cG3!45] (v5) -- node[font=\tiny,fill=white,inner sep=0.4pt]{1} (v6);
\draw[cG1,thick,dashed] (v5) -- (v6);
\draw[line width=2.6pt,cG3!45] (v3) -- node[font=\tiny,fill=white,inner sep=0.4pt]{14} (v5);
\draw[line width=2.6pt,cG3!45] (v0) -- node[font=\tiny,fill=white,inner sep=0.4pt]{11} (v6);
\draw[cG1,thick,dashed] (v0) -- (v6);
\draw[line width=2.6pt,cG3!45] (v6) -- node[font=\tiny,fill=white,inner sep=0.4pt]{10} (v7);
\draw[cG1,thick,dashed] (v6) -- (v7);
\draw[black!30] (v1) -- node[font=\tiny,fill=white,inner sep=0.4pt]{13} (v3);
\draw[cG1,thick,dashed] (v1) -- (v3);
\draw[line width=2.6pt,cG3!45] (v2) -- node[font=\tiny,fill=white,inner sep=0.4pt]{11} (v5);
\draw[cG1,thick,dashed] (v2) -- (v5);
\draw[black!30] (v0) -- node[font=\tiny,fill=white,inner sep=0.4pt]{13} (v8);
\end{tikzpicture}
\caption{G3: cost 58, OPT 48, ratio 1.208.}\end{subfigure}
\caption{Worst instances found by the search, one per strategy. Filled vertices are terminals. The wide light band marks an optimal tree (from the DP); the dashed line is the tree the greedy strategy returns. Numbers on edges are weights.}\label{fig:gadgets}\end{figure}


\subsection{Optimality gaps on the benchmark}\label{sec:gaps}
\Cref{tab:gaps} compares the three strategies on every instance with a known optimum (published for SteinLib, from the dynamic
program for the generated sets). The gap is \((\mathrm{cost}-\OPT)/\OPT\).

\begin{table}[!htbp]\centering\small
\begin{tabular}{@{}lrrrrrrrrrrrrr@{}}\toprule
 & & \multicolumn{4}{c}{G1 KMB} & \multicolumn{4}{c}{G2 SPH (best root)} & \multicolumn{4}{c}{G3 Zelikovsky}\\
\cmidrule(lr){3-6}\cmidrule(lr){7-10}\cmidrule(l){11-14}
Set & \# & mean & max & opt & ms & mean & max & opt & ms & mean & max & opt & ms\\\midrule
B & 18 & 2.38 & 6.30 & 7 & 1.0 & \textbf{0.18} & 1.64 & 15 & 11.3 & 0.65 & 4.72 & 12 & 1.5\\
C & 20 & 5.36 & 15.62 & 1 & 26.1 & \textbf{1.99} & 9.09 & 5 & 1703.0 & 2.33 & 9.09 & 3 & 153.9\\
P6E & 15 & 2.77 & 8.96 & 2 & 1.5 & 0.36 & 3.49 & 10 & 39.1 & \textbf{0.09} & 1.27 & 13 & 1.9\\
I080 & 100 & 18.51 & 41.03 & 1 & 0.6 & 10.73 & 41.03 & 13 & 6.5 & \textbf{8.14} & 20.48 & 3 & 0.9\\
Grid & 27 & 2.89 & 9.41 & 6 & 0.5 & \textbf{0.54} & 2.68 & 17 & 2.4 & 0.57 & 2.72 & 17 & 0.6\\
Random & 6 & 7.95 & 14.27 & 0 & 4.5 & \textbf{1.03} & 2.04 & 1 & 22.1 & 2.42 & 5.30 & 0 & 4.6\\
Adversarial & 13 & 50.48 & 74.17 & 1 & 0.1 & 48.78 & 74.17 & 2 & 0.2 & \textbf{1.60} & 20.83 & 12 & 0.1\\
\midrule All & 199 & 14.19 & 74.17 & 18 & 3.3 & 8.93 & 74.17 & 63 & 179.4 & 4.65 & 20.83 & 60 & 16.4\\
\bottomrule\end{tabular}
\caption{Gap to the optimum in percent (mean and maximum), number of instances solved optimally, and mean time in milliseconds. Bold: smallest mean gap in the row.}\label{tab:gaps}
\end{table}


\noindent First, all three are far better on the benchmarks than their worst case: the largest gap on a SteinLib instance is
\val{gap.max.steinlib}\%, nowhere near a ratio of \(11/6\) or \(2\); the gaps of 50--74\% appear only on the adversarial set, which was built to
cause them. Second, no strategy wins everywhere. G3 has the smallest mean gap overall (\val{g3.meangap}\%) because it is the only one that
stays close to the optimum on the adversarial instances and on I080, whose incidence weights punish pairwise path joins. G2 has the
smallest gap on B, C, grids and random graphs. Third, G2 pays for that with time: it runs the heuristic once per terminal, which makes it
the slowest greedy strategy by an order of magnitude (\val{time.C.G2}\,s per C instance on average, against \val{time.C.G3}\,s for G3 and \val{time.C.G1}\,s for G1). With a single root
it is much weaker, and its worst root is more expensive than G1 on \val{sph.worse_than_g1} of \val{gap.count} instances
(\cref{tab:sphroots}). The choice of root is part of the greedy decision, and a bad one costs as much as the difference between G1 and G2.

\begin{table}[H]\centering\small
\begin{tabular}{@{}lrrrrr@{}}\toprule
Set & G1 & G2, first terminal & G2, worst root & G2, best root & worst \(>\) G1\\\midrule
B & 2.38 & 1.57 & 2.65 & 0.18 & 5/18\\
C & 5.36 & 2.97 & 4.57 & 1.99 & 7/20\\
P6E & 2.77 & 1.17 & 2.53 & 0.36 & 5/15\\
I080 & 18.51 & 14.62 & 17.28 & 10.73 & 22/100\\
Grid & 2.89 & 1.37 & 2.80 & 0.54 & 6/27\\
Random & 7.95 & 3.23 & 4.10 & 1.03 & 2/6\\
Adversarial & 50.48 & 48.78 & 50.38 & 48.78 & 1/13\\
\bottomrule\end{tabular}
\caption{Mean gap (\%) of G2 depending on the root terminal, next to G1. The last column counts instances where the worst root gives a more expensive tree than G1.}\label{tab:sphroots}
\end{table}


% ------------------------------------------------------------------ 5
\section{Exact paradigm 1: dynamic programming over terminal subsets}\label{sec:dp}

\subsection{State and recurrence}
Fix a root terminal \(r\) and let \(Q=T\setminus\{r\}\), \(q=|Q|=k-1\). The state is a pair \((S,v)\) with
\(\emptyset\neq S\subseteq Q\) and \(v\in V\):
\[\mathrm{dp}[S][v]=\text{cost of a cheapest tree in }G\text{ that contains }S\cup\{v\}.\]
The answer is \(\mathrm{dp}[Q][r]\), a cheapest tree containing \(Q\cup\{r\}=T\). The recurrence has a base case and two
steps:
\begin{align}
\mathrm{dp}[\{t\}][v]&=d(t,v), \label{eq:base}\\
\mathrm{mrg}[S][u]&=\min_{\emptyset\neq S_1\subsetneq S}\ \mathrm{dp}[S_1][u]+\mathrm{dp}[S\setminus S_1][u] \qquad(|S|\ge2), \label{eq:merge}\\
\mathrm{dp}[S][v]&=\min_{u\in V}\ \mathrm{mrg}[S][u]+d(u,v) \qquad(|S|\ge2). \label{eq:grow}
\end{align}
The merge step \eqref{eq:merge} joins two trees that meet at \(u\); the grow step \eqref{eq:grow} hangs that junction off
\(v\) by a shortest path.

\begin{lemma}[optimal substructure]\label{lem:dw}
\eqref{eq:base}--\eqref{eq:grow} compute \(\mathrm{dp}[S][v]\) correctly for every state.
\end{lemma}
\begin{proof}
For \(|S|=1\) the cheapest tree containing \(t\) and \(v\) is a shortest path. Let \(|S|\ge2\).

(\(\ge\)) Take an optimal tree \(A\) for \(S\cup\{v\}\); we may assume every leaf of \(A\) is in \(S\cup\{v\}\). Walk from
\(v\) along \(A\) while the current vertex has exactly one unvisited neighbour in \(A\) and is not in \(S\) (if \(v\) itself is a
branch point or in \(S\), the walk stops at once). It stops at a vertex \(u\) that is a branch point or a terminal of \(S\). The walk
is a \(v\)--\(u\) path of cost at least \(d(u,v)\). The rest of \(A\) is a tree containing \(S\cup\{u\}\) in which either \(u\) has
two or more subtrees, or \(u\in S\). In the first case split the subtrees of \(u\) into two non-empty groups; each group contains a
terminal of \(S\) (every leaf is in \(S\)), giving trees for \(S_1\cup\{u\}\) and \((S\setminus S_1)\cup\{u\}\). In the second
case take \(S_1=\{u\}\), whose tree is the single vertex \(u\) with \(\mathrm{dp}[\{u\}][u]=0\). Either way the cost of \(A\) is at least
\(d(v,u)+\mathrm{dp}[S_1][u]+\mathrm{dp}[S\setminus S_1][u]\ge\) the right-hand side of \eqref{eq:grow}.

(\(\le\)) For any \(u\) and split \(S_1\), the union of a \(v\)--\(u\) shortest path and the two optimal subtrees is a connected
subgraph containing \(S\cup\{v\}\); it contains a tree of no greater cost.
\end{proof}

\noindent Subsets are processed in increasing numerical order of their bitmask, so both halves of every split are finished before \(S\)
is merged.

\subsection{The grow step is a Dijkstra run}
Equation \eqref{eq:grow} is a shortest-path problem with \emph{initial labels}: put \(\mathrm{mrg}[S][u]\) on every vertex and
run Dijkstra from all of them at once. The extracted label of \(v\) is \(\min_u \mathrm{mrg}[S][u]+d(u,v)\), by the usual Dijkstra
argument applied to a super-source joined to each \(u\) by an edge of weight \(\mathrm{mrg}[S][u]\). This is the
Erickson--Monma--Veinott form of the Dreyfus--Wagner algorithm. The original paper precomputes all-pairs distances and
evaluates \eqref{eq:grow} as a full \(O(n^2)\) minimisation per subset. With a Fibonacci heap one layer costs \(n\) \textsc{Insert},
\(n\) \textsc{Extract-Min} and at most \(m\) \textsc{Decrease-Key}, that is \(O(m+n\log n)\) by \cref{thm:sequence}.

\begin{algorithm}[!htbp]
\For{\(S=1\) \KwTo \(2^{q}-1\)}{
  \eIf{\(S=\{t_i\}\)}{\(\mathrm{lab}[\cdot]\gets\infty\); \(\mathrm{lab}[t_i]\gets0\)\;}{
    \(\mathrm{lab}[\cdot]\gets\infty\); \(\ell\gets\) lowest bit of \(S\)\;
    \ForEach(\tcp*[f]{submask enumeration}){\(S_1\subsetneq S\) with \(\ell\in S_1\)}{
      \lFor{\(v\in V\)}{\(\mathrm{lab}[v]\gets\min(\mathrm{lab}[v],\ \mathrm{dp}[S_1][v]+\mathrm{dp}[S\setminus S_1][v])\)}
    }
  }
  multi-source Dijkstra on the Fibonacci heap with initial labels \(\mathrm{lab}\)\;
  \(\mathrm{dp}[S][\cdot]\gets\mathrm{lab}\); \(\mathrm{pred}[S][\cdot]\gets\) last edge of each Dijkstra path (or \(-1\))\;
}
\Return \(\mathrm{dp}[Q][r]\) and the tree read back from \(\mathrm{pred}\)\;
\caption{DP (\texttt{src/dp\_dreyfus\_wagner.cpp}). The inner loop over \(v\) is contiguous in memory and vectorises.}
\end{algorithm}

\subsection{Complexity and memory}
\begin{theorem}\label{thm:dpcost}
The DP runs in \(O\big(3^{q}n+2^{q}(m+n\log n)\big)\) time and \(2^{q}n\) words of memory, \(q=k-1\).
\end{theorem}
\begin{proof}
Requiring the lowest bit of \(S\) to be in \(S_1\) visits each unordered split \(\{S_1,S\setminus S_1\}\) exactly once:
\(2^{|S|-1}-1\) splits for a set of size \(|S|\ge2\). Summing over all subsets,
\[\sum_{s=2}^{q}\binom{q}{s}\big(2^{s-1}-1\big)=\frac{3^{q}+1}{2}-2^{q}\]
merges, each a loop over \(n\) vertices. (For \(k=9\) that is 3025 merges, and the program reports exactly 3025.) There are
\(2^{q}-1\) Dijkstra runs at \(O(m+n\log n)\) each, and the table has \(2^{q}\) rows of \(n\) entries.
\end{proof}

\noindent Fixing the root saves a factor 3 in time and 2 in memory over a formulation that ranges over subsets of all \(k\)
terminals. Each state stores an 8-byte cost and a 4-byte predecessor edge, so the table needs \(12\cdot2^{q}n\) bytes:
\val{dp.mem.k17}\,MB for \(k=17\), \(n=100\) and \val{dp.mem.k20}\,MB for \(k=20\), \(n=80\). The solver refuses an instance whose
table would not fit in the memory limit (8\,GB by default) instead of swapping, and it is run with \(k\le20\).

\paragraph{Reading the tree back.} From \((Q,r)\) the solver follows \(\mathrm{pred}\) inside the same layer. When the
predecessor is \(-1\) the label came from the merge step, and it looks for a split whose two halves add up to the label exactly; one
exists by \eqref{eq:merge}. The traced edges form a connected subgraph containing \(T\) whose cost, counting shared edges once, is at
most \(\mathrm{dp}[Q][r]\). \Cref{lem:cleanup} turns it into a tree of cost at most \(\OPT\), hence exactly \(\OPT\). The program checks
that the tree cost equals the DP value and the validator checks the tree.

\subsection{Results}
\begin{table}[!htbp]\centering\small
\begin{tabular}{@{}lrrrrr@{}}\toprule
Set & instances & with \(k\le20\) & solved & equal to OPT & max time (s)\\\midrule
B & 18 & 10 & 10 & 10 & 27.19\\
C & 20 & 8 & 8 & 8 & 0.24\\
P6E & 15 & 10 & 9 & 9 & 155.13\\
I080 & 100 & 100 & 100 & 100 & 159.31\\
Grid & 27 & 27 & 27 & 27 & 124.69\\
Random & 6 & 6 & 6 & 6 & 1.12\\
Adversarial & 13 & 13 & 13 & 13 & 0.01\\
\bottomrule\end{tabular}
\caption{DP results. Instances with \(k>20\) are not attempted (the table would need more than \(12\cdot2^{20}n\) bytes and \(3^{20}n\) merges). Every solved instance matches the known optimum.}\label{tab:dp}
\end{table}
\begin{figure}[!htbp]\centering
\begin{tikzpicture}
\begin{semilogyaxis}[width=0.48\linewidth,height=5.4cm,xlabel={terminals \(k\)},ylabel={time (s)},
  legend pos=north west,legend style={font=\scriptsize,draw=none},tick label style={font=\footnotesize},label style={font=\small},
  grid=major,grid style={black!8},xtick={4,6,...,20}]
\addplot[cDP,thick,mark=*,mark size=1.6pt] table[x=k,y=grid_s,col sep=comma]{generated/dp_vs_k.csv};\addlegendentry{grid, \(n=100\)}
\addplot[cG3,thick,mark=square*,mark size=1.6pt] table[x=k,y=i080_s,col sep=comma]{generated/dp_vs_k.csv};\addlegendentry{I080, \(n=80\)}
\addplot[black!50,dashed,domain=4:20,samples=17] {9.926448e-10*100*3^(x-1)};\addlegendentry{\(c\cdot3^{k-1}n\)}
\end{semilogyaxis}
\end{tikzpicture}\hfill
\begin{tikzpicture}
\begin{loglogaxis}[width=0.48\linewidth,height=5.4cm,xlabel={vertices \(n\) (\(k=10\))},ylabel={time (s)},
  legend pos=north west,legend style={font=\scriptsize,draw=none},tick label style={font=\footnotesize},label style={font=\small},
  grid=major,grid style={black!8}]
\addplot[cDP,thick,mark=*,mark size=1.6pt] table[x=n,y=dp_s,col sep=comma]{generated/dp_vs_n.csv};\addlegendentry{DP}
\addplot[cBB,thick,mark=triangle*,mark size=1.8pt] table[x=n,y=bnb_s,col sep=comma]{generated/dp_vs_n.csv};\addlegendentry{B\&B}
\addplot[cG1,thick,mark=o,mark size=1.6pt] table[x=n,y=g1_s,col sep=comma]{generated/dp_vs_n.csv};\addlegendentry{G1}
\addplot[cG3,thick,mark=square,mark size=1.6pt] table[x=n,y=g3_s,col sep=comma]{generated/dp_vs_n.csv};\addlegendentry{G3}
\end{loglogaxis}
\end{tikzpicture}
\caption{Left: DP time against \(k\) (mean of the instances per \(k\)); the dashed line is \(3^{k-1}n\) scaled to the \(k=20\) grids,
and the measured curve follows it. Right: time against \(n\) on sparse random graphs with \(k=10\).}\label{fig:dpscale}
\end{figure}


% ------------------------------------------------------------------ 6
\section{Exact paradigm 2: branch-and-bound}\label{sec:bnb}

The DP is exponential in \(k\) and useless beyond \(k\approx20\). Branch-and-bound is exponential in the number of branching
decisions instead, so it helps when a good lower bound settles most decisions before any branching happens. Our design follows the
structure of the strongest exact Steiner solvers (Koch and Martin 1998; Polzin and Vahdati Daneshmand 2001), cut down to what can be
written without an LP solver: exact graph reductions, a dual ascent lower bound, a heuristic upper bound, reduced-cost
elimination, and best-first branching on Steiner vertices.

\subsection{Search space and branching}
A search node fixes some non-terminals as \emph{included} (they become terminals) and some vertices and edges as \emph{excluded}
(they are deleted). Branching on a free non-terminal \(v\) creates two children, \(v\) included and \(v\) excluded. Every Steiner tree
either contains \(v\) or does not, so the two children together keep every feasible tree. Each branching fixes one vertex, so every
root-to-leaf path has at most \(n-k\) branchings, and the search tree is finite with at most \(2^{n-k+1}\) nodes. We branch on the free
non-terminal with the smallest \(d'(r,v)+d'(v,T)\) (defined in \cref{sec:elim}): the vertex that the dual solution considers cheapest
to use, which is most likely in an optimal tree and therefore tightens the bound quickly in the excluding child.

Open nodes wait in a Fibonacci heap keyed by their lower bound, and the node with the smallest bound is expanded next
(best-first search). The minimum key is a global lower bound on \(\OPT\): as soon as it reaches the incumbent cost \(U\) the search
stops with a proof of optimality.

\subsection{Exact reductions}
Before the search the graph is reduced. All four tests keep at least one optimal tree.

\begin{lemma}\label{lem:reductions}
Let \(k\ge2\). (R1) A non-terminal of degree at most 1 lies in no optimal tree with only terminal leaves; delete it.
(R2) If a terminal \(t\) has a single edge \(\{t,u\}\), that edge is in every feasible tree; add \(w(t,u)\) to a fixed cost, delete \(t\)
and make \(u\) a terminal. (R3) If a non-terminal \(v\) has exactly two edges \(\{a,v\},\{v,b\}\), replace them by one edge
\(\{a,b\}\) of weight \(w(a,v)+w(v,b)\), keeping the cheaper one if \(\{a,b\}\) already exists. (R4) Every edge with
\(w(u,v)>d(u,v)\) can be deleted, all at once.
\end{lemma}
\begin{proof}
(R1) A non-terminal leaf can be pruned without raising the cost. (R2) \(t\) must be connected to another terminal, and its only
edge is the only way. (R3) In a tree with only terminal leaves, \(v\) either does not appear or has degree 2 and uses both edges, which
is exactly the new edge; if a cheaper parallel edge exists, any tree using the path can switch to it. (R4) If a tree uses
\(\{u,v\}\), replace that edge by a shortest \(u\)--\(v\) path. The result is connected, spans \(T\) and is strictly cheaper, so the
edge is in no optimal tree. Edges that are strictly longer than the distance between their ends lie on no shortest path at all, so
deleting them changes no distance and the argument holds for all of them together.
\end{proof}
\noindent R1--R3 run until nothing changes, then R4 (one Dijkstra per vertex), and the whole cycle repeats while anything changes.
Each reduced edge remembers the list of original edges it stands for, so the final tree is mapped back to the input graph and checked
there.

\subsection{Lower bound: dual ascent on the cut formulation}
Orient the problem from the root: replace each edge by two opposite arcs of the same weight. A Steiner tree oriented away from \(r\) is
an arborescence \(A\), and \(x=\chi^{A}\) satisfies
\[\min\ \sum_{a}c_a x_a\quad\text{s.t.}\quad x(\delta^-(W))\ge1\ \ \text{for all }W\subseteq V,\ r\notin W,\ W\cap T\neq\emptyset,\qquad x\ge0, \tag{P}\]
where \(\delta^-(W)\) is the set of arcs entering \(W\). The dual is
\[\max\ \sum_W y_W\quad\text{s.t.}\quad \sum_{W:\,a\in\delta^-(W)}y_W\le c_a\ \ \text{for all arcs }a,\qquad y\ge0. \tag{D}\]

\begin{lemma}[admissibility]\label{lem:lb}
For a dual-feasible \(y\) with reduced costs \(c'_a=c_a-\sum_{W:a\in\delta^-(W)}y_W\ge0\), every Steiner arborescence \(A\) satisfies
\(c(A)\ge\sum_W y_W+c'(A)\). In particular \(\mathrm{LB}=\sum_Wy_W\le\OPT\).
\end{lemma}
\begin{proof}
\(c(A)=\sum_{a\in A}c'_a+\sum_{a\in A}\sum_{W:a\in\delta^-(W)}y_W=c'(A)+\sum_W y_W\,|A\cap\delta^-(W)|\), and
\(|A\cap\delta^-(W)|\ge1\) because \(A\) contains a path from \(r\notin W\) to a terminal in \(W\).
\end{proof}

Wong's dual ascent builds such a \(y\) greedily. For an active terminal \(t\), let \(W(t)\) be the set of vertices that reach \(t\)
through arcs of reduced cost 0. While \(r\notin W(t)\), raise \(y_{W(t)}\) by the smallest reduced cost on an arc entering \(W(t)\);
that arc becomes saturated (\(c'=0\)), its tail joins \(W(t)\), and the bound grows by the same amount. Reduced costs never go
negative, so \(y\) stays feasible and \cref{lem:lb} applies at every moment. Our implementation grows one component at a time and keeps
the arcs entering \(W(t)\) in a Fibonacci heap. Raising \(y\) by \(\delta\) lowers every entering arc by \(\delta\), so instead of
touching them it adds \(\delta\) to an offset \(Y\) and keys each arc by its reduced cost plus the value of \(Y\) when it entered the heap:
the heap minimum is always the next arc to saturate. A component stops growing when it reaches \(r\) or another active terminal
(that terminal will connect both). Each growth step saturates one arc and arcs never unsaturate, so a call makes at most \(2m\) steps
per terminal, \(O(km\log m)\) in total. At the root node the ascent is repeated from up to 16 different root terminals and the best bound
is kept.

\subsection{Upper bound and reduced-cost elimination}\label{sec:elim}
At every node the shortest-path heuristic (\cref{sec:greedy}, G2) runs on the node graph and on the subgraph of arcs with reduced cost
0. That subgraph connects every terminal to \(r\), and trees inside it tend to be cheap because the dual has already paid for those
arcs. At the root it runs from up to 64 start terminals. The best tree is then improved by local search on its vertex set (try deleting
each Steiner vertex and adding each neighbouring vertex, recompute the MST, keep improvements). This search runs at the root, on every
node whose heuristic beats the incumbent, and on every 64th node.

\begin{lemma}[reduced-cost elimination]\label{lem:elim}
Let \(d'(x,y)\) be the shortest-path distance under the reduced costs \(c'\), \(d'(v,T)=\min_{t\in T}d'(v,t)\), and \(U\) the cost of
the incumbent. A non-terminal \(v\) with \(\mathrm{LB}+d'(r,v)+d'(v,T)\ge U\) lies in no tree cheaper than \(U\). An arc
\((u,v)\) with \(\mathrm{LB}+d'(r,u)+c'_{uv}+d'(v,T)\ge U\) lies in no arborescence cheaper than \(U\).
\end{lemma}
\begin{proof}
Let \(A\) be an arborescence with only terminal leaves that contains \(v\). It contains a path from \(r\) to \(v\) and, below \(v\), a path
from \(v\) to a terminal; the two share no arc. Since \(c'\ge0\), \(c'(A)\ge d'(r,v)+d'(v,T)\), and \cref{lem:lb} gives
\(c(A)\ge\mathrm{LB}+d'(r,v)+d'(v,T)\ge U\). The arc case is the same with the path \(r\to u\to v\to T\). An edge is deleted when
both of its arcs fail the test.
\end{proof}
\noindent Both \(d'\) are one Dijkstra run on the heap (from \(r\), and backwards from all terminals). The deletions are valid for the
whole subtree of the node, because \(U\) can only decrease, so children inherit them.

\begin{theorem}\label{thm:bnb}
If the search ends before the limits, the incumbent is optimal.
\end{theorem}
\begin{proof}
Invariant: either the incumbent is optimal, or some optimal tree is feasible for an open node. It holds at the start, and it is kept by
branching (the children partition the trees of the parent), by the reductions (\cref{lem:reductions}) and by the eliminations
(\cref{lem:elim}), which only remove trees that cost at least \(U\). A node is discarded only when \(\mathrm{LB}\ge U\), and by
\cref{lem:lb} it then holds no tree cheaper than \(U\). When the heap is empty, or its smallest key is at least \(U\), no open node can
hold a cheaper tree, so the incumbent is optimal. Weights are integers, so \(\mathrm{LB}\) is an integer and the comparison is exact.
\end{proof}

\paragraph{Complexity.} One node costs \(O(m)\) to build, \(O(km\log m)\) for the dual ascent, \(O(k(m+n\log n))\) for each heuristic run,
\(O(m+n\log n)\) for the two reduced-cost Dijkstra runs and \(O(n(m+n\log n))\) per local search pass when it runs. The number of nodes is
at most \(2^{n-k+1}\) in the worst case, so the worst-case running time is exponential, as it must be unless P\,=\,NP. The
measurements below show how far below that bound the search stays.

\subsection{Results}
\begin{table}[H]\centering\small\setlength{\tabcolsep}{4pt}
\begin{tabular}{@{}lrrrrrrrrr@{}}\toprule
Set & \# & solved & \(|V|\) left & \(|E|\) left & root gap \% & nodes: median & max & mean s & max s\\\midrule
B & 18 & 18 & 60\% & 67\% & 0.05 & 1 & 3 & 0.01 & 0.05\\
C & 20 & 19 & 79\% & 65\% & 0.61 & 3 & 23{,}088 & 17.17 & 300.02\\
P6E & 15 & 15 & 97\% & 98\% & 0.22 & 1 & 13{,}707 & 1.10 & 13.06\\
I080 & 100 & 100 & 86\% & 91\% & 2.31 & 10 & 835 & 0.07 & 0.53\\
Grid & 27 & 27 & 90\% & 89\% & 0.06 & 1 & 8 & 0.02 & 0.04\\
Random & 6 & 6 & 75\% & 77\% & 0.42 & 28 & 64 & 1.95 & 8.73\\
Adversarial & 13 & 13 & 93\% & 93\% & 0.00 & 1 & 1 & 0.00 & 0.00\\
\bottomrule\end{tabular}
\caption{Branch-and-bound with a 300\,s limit. \(|V|\), \(|E|\) left: share of the graph that survives the reductions. Root gap: \((\OPT-\mathrm{LB}_{\text{root}})/\OPT\) for the dual ascent bound at the root.}\label{tab:bnb}
\end{table}
\noindent Instances not solved within the limit, with the bounds proved when the search stopped:
\begin{center}\small\begin{tabular}{@{}lrrrrr@{}}\toprule
Instance & \(k\) & lower bound & best tree & OPT & nodes\\\midrule
c18 & 83 & 110 & 117 & 113 & 23{,}088\\
\bottomrule\end{tabular}\end{center}
\begin{figure}[H]\centering
\begin{tikzpicture}
\begin{loglogaxis}[width=0.42\linewidth,height=5.6cm,xlabel={DP time (s)},ylabel={B\&B time (s)},
  colorbar,colormap={kk}{color=(cG3) color=(cG1)},point meta=explicit,colorbar style={ylabel={\(k\)},font=\footnotesize,width=2.5mm},
  tick label style={font=\footnotesize},label style={font=\small},grid=major,grid style={black!8}]
\addplot[scatter,only marks,mark=*,mark size=1.5pt,fill opacity=0.7,draw opacity=0] table[x=dp,y=bb,meta=k,col sep=comma]{generated/dp_vs_bb.csv};
\addplot[black!50,dashed,domain=1e-5:200] {x};
\end{loglogaxis}
\end{tikzpicture}
\caption{The two exact methods on the 173 instances both solve. Below the diagonal B\&B is faster; it is faster on 97 of them.
Colour is the number of terminals: the DP time is set by \(k\), the B\&B time by how much of the graph the bound settles.}\label{fig:dpbb}
\end{figure}


% ------------------------------------------------------------------ 7
\section{Experimental setup and overall results}\label{sec:experiments}

\paragraph{Machine and build.} \val{env.cpu}, \val{env.cores} logical cores, \val{env.mem}\,GB RAM, \val{env.os}; \val{env.gxx},
flags \texttt{-std=c++17 -O2 -march=native}. The main table ran six solver processes in parallel
(\texttt{scripts/run\_bench.sh}); the heap comparison of \cref{sec:ablation} ran one process at a time and keeps the best of three
runs, because it compares times. The exact solvers have a 300\,s limit per instance.

\paragraph{Instances.} SteinLib sets B (18 sparse random graphs, \(n\le100\)), C (20 graphs with \(n=500\), up to \(12{,}500\) edges and
\(k\) up to 250), P6E (15 Euclidean-weight graphs, \(n=100\)) and I080 (100 graphs with \(n=80\) and Duin's incidence weights, which
are designed to make reduction tests ineffective), all with published optima. The generated sets are the adversarial instances of \cref{sec:counter}, 27 grids
(\(10\times10\), \(k=4,6,\dots,20\), three each) and six sparse random graphs with \(n\) from 100 to 3200 and \(k=10\). In total
\val{inst.total} instances and \val{rows.total} solver runs.

\paragraph{Correctness checks.} Every tree any solver returns passes through the validator in \texttt{include/graph.hpp},
which is separate from all solvers: edge ids valid and distinct, no cycle, connected, all terminals spanned, and the cost recomputed from
the input. No run failed this check. \texttt{scripts/make\_tables.py} also refuses to build the report if any greedy cost is below the
known optimum, or if an exact solver that claims optimality differs from it. All \val{exact.checked} exact SteinLib results agree with the
published optima, and on the generated sets the DP and the branch-and-bound, two unrelated algorithms,
agree with each other. Finally, \texttt{scripts/fuzz.py} runs all five solvers on 300 random graphs with up to 14 vertices,
zero-weight edges, ties and \(k\) from 1 to \(n\); the two exact methods agree on all of them and no greedy tree is invalid or below
the optimum.

\begin{table}[!htbp]\centering\small
\begin{tabular}{@{}llrrrr@{}}\toprule
Solver & Paradigm & runs & optimal & mean gap (\%) & total time (s)\\\midrule
G1 & greedy & 199 & 18 & 14.195 & 0.7\\
G2 & greedy & 199 & 63 & 8.927 & 35.7\\
G3 & greedy & 199 & 60 & 4.646 & 3.3\\
DP & dynamic programming & 173 & 173 & 0.000 & 4617.3\\
B\&B & branch-and-bound & 199 & 198 & 0.018 & 379.3\\
\bottomrule\end{tabular}
\caption{All instances with a known optimum (199). \emph{runs}: runs that returned a tree (the DP skips \(k>20\)). For the exact methods a non-zero mean gap comes only from instances stopped at the time limit.}\label{tab:overview}
\end{table}


\subsection{Fibonacci heap against binary heap}\label{sec:ablation}
\begin{table}[H]\centering\small
\begin{tabular}{@{}llrrrrr@{}}\toprule
Solver & graphs & instances & Fibonacci (s) & binary (s) & Fib.\,/\,bin. & Fib. faster\\\midrule
G1 & sparse, \(m/n<10\) & 77 & 0.024 & 0.010 & 2.33 & 0/77\\
G1 & dense, \(m/n\ge10\) & 7 & 0.003 & 0.002 & 1.51 & 0/7\\
G2 & sparse, \(m/n<10\) & 77 & 0.146 & 0.056 & 2.60 & 0/77\\
G2 & dense, \(m/n\ge10\) & 7 & 0.010 & 0.006 & 1.58 & 0/7\\
G3 & sparse, \(m/n<10\) & 77 & 0.027 & 0.013 & 2.05 & 0/77\\
G3 & dense, \(m/n\ge10\) & 7 & 0.003 & 0.002 & 1.50 & 0/7\\
DP & sparse, \(m/n<10\) & 59 & 3.573 & 2.022 & 1.77 & 0/59\\
DP & dense, \(m/n\ge10\) & 7 & 0.072 & 0.055 & 1.31 & 0/7\\
\bottomrule\end{tabular}
\caption{Total time with each heap (best of three serial runs per instance). Sparse: SteinLib B, I080 (\(k\le16\) for the DP), grids and random graphs; dense: SteinLib c11, c12, c16, c17 (\(m/n=5\) to \(25\)) with c01, c02, c06, c07 as sparse C graphs. The DP returns the same optimum with both heaps; the greedy strategies can break ties between equal keys differently.}\label{tab:ablation}
\end{table}

\noindent The binary heap is faster on every instance, by a factor of \val{abl.DP.sparse.ratio} for the DP on sparse graphs and
\val{abl.G2.sparse.ratio} for G2. This does not contradict \cref{thm:sequence}: the theorem bounds the number of pointer
operations, not the time each one takes. The binary heap keeps its keys in one array that stays in cache, while every Fibonacci
node is a separate allocation with four pointers. On these graphs the asymptotic advantage has little to work with, because a
Dijkstra run on a graph with \(m/n<10\) does fewer decrease-keys than extract-mins (\cref{tab:heapcounts}), and each binary-heap
decrease-key moves a key only a few levels.

The gap does narrow where the theory says it should. On the dense graphs the DP ratio falls to \val{abl.DP.dense.ratio}, and
on \texttt{c17} (\(m/n=25\)) the run makes \val{abl.dp.c17.decr} decrease-keys against \val{abl.dp.c17.extr} extract-mins and
takes \val{abl.dp.c17.fib}\,ms with the Fibonacci heap against \val{abl.dp.c17.bin}\,ms with the binary heap. We report the
Fibonacci heap as the structure that gives the proved bounds, not as the faster choice on SteinLib-sized inputs.


\section{Limitations}\label{sec:limits}
\begin{itemize}
\item The DP is limited to \(k\le20\) by time and memory (\cref{thm:dpcost}). A meet-in-the-middle or the Nederlof
inclusion--exclusion variant would trade time for polynomial space, but would no longer use the heap.
\item The branch-and-bound has no LP solver, so its bound is the dual ascent value, which is weaker than the LP optimum on some dense
graphs. The instances left open in \cref{tab:bnb} are those where the root gap stays large; LP-based branch-and-cut solvers such as
SCIP-Jack are the standard tool for them.
\item The search for adversarial instances only covers graphs with at most 9 vertices. It shows that each greedy strategy can be
suboptimal, not how bad it can be: the ratios found (at most \val{adv.g1.ratio}) are far from the proven bounds.
\item Times are wall-clock times on one machine, measured inside each process from after the input is read until the solver returns its tree;
runs in the main table shared the machine with five others.
\end{itemize}


\begin{thebibliography}{99}\small
\bibitem{ft87} M.~L. Fredman and R.~E. Tarjan. Fibonacci heaps and their uses in improved network optimization algorithms. \emph{J. ACM} 34(3):596--615, 1987.
\bibitem{clrs} T.~H. Cormen, C.~E. Leiserson, R.~L. Rivest, C.~Stein. \emph{Introduction to Algorithms}, 3rd ed., ch.~17 and 19. MIT Press, 2009.
\bibitem{tarjan85} R.~E. Tarjan. Amortized computational complexity. \emph{SIAM J. Algebraic Discrete Methods} 6(2):306--318, 1985.
\bibitem{karp72} R.~M. Karp. Reducibility among combinatorial problems. In \emph{Complexity of Computer Computations}, 85--103, 1972.
\bibitem{chlebik} M.~Chleb\'ik and J.~Chleb\'ikov\'a. The Steiner tree problem on graphs: inapproximability results. \emph{Theor. Comput. Sci.} 406(3):207--214, 2008.
\bibitem{dw71} S.~E. Dreyfus and R.~A. Wagner. The Steiner problem in graphs. \emph{Networks} 1(3):195--207, 1971.
\bibitem{emv87} R.~E. Erickson, C.~L. Monma, A.~F. Veinott. Send-and-split method for minimum-concave-cost network flows. \emph{Math. Oper. Res.} 12(4):634--664, 1987.
\bibitem{kmb81} L.~Kou, G.~Markowsky, L.~Berman. A fast algorithm for Steiner trees. \emph{Acta Informatica} 15:141--145, 1981.
\bibitem{tm80} H.~Takahashi and A.~Matsuyama. An approximate solution for the Steiner problem in graphs. \emph{Math. Japonica} 24(6):573--577, 1980.
\bibitem{mehlhorn} K.~Mehlhorn. A faster approximation algorithm for the Steiner problem in graphs. \emph{Inf. Process. Lett.} 27(3):125--128, 1988.
\bibitem{zel93} A.~Z. Zelikovsky. An 11/6-approximation algorithm for the network Steiner problem. \emph{Algorithmica} 9(5):463--470, 1993.
\bibitem{wong84} R.~T. Wong. A dual ascent approach for Steiner tree problems on a directed graph. \emph{Math. Programming} 28:271--287, 1984.
\bibitem{km98} T.~Koch and A.~Martin. Solving Steiner tree problems in graphs to optimality. \emph{Networks} 32(3):207--232, 1998.
\bibitem{pd01} T.~Polzin and S.~Vahdati Daneshmand. Improved algorithms for the Steiner problem in networks. \emph{Discrete Appl. Math.} 112:263--300, 2001.
\bibitem{steinlib} T.~Koch, A.~Martin, S.~Vo\ss. SteinLib: an updated library on Steiner tree problems in graphs. ZIB report 00-37, 2000. \url{https://steinlib.zib.de}
\bibitem{beasley} J.~E. Beasley. An algorithm for the Steiner problem in graphs. \emph{Networks} 14:147--159, 1984.
\bibitem{bgrs} J.~Byrka, F.~Grandoni, T.~Rothvo\ss, L.~Sanit\`a. Steiner tree approximation via iterative randomized rounding. \emph{J. ACM} 60(1):6, 2013.
\end{thebibliography}

\appendix
\section{Reproducing the results}\label{app:repro}
\begin{verbatim}
make                              # the five solvers, in build/
make test                         # heap tests, amortized check
python3 scripts/fuzz.py           # 300 random graphs, all solvers
python3 scripts/gen_instances.py  # generated instances
bash scripts/run_bench.sh         # results/*.csv
python3 scripts/make_tables.py    # report/generated/*
make report                       # report/report.pdf
\end{verbatim}
Every solver reads a SteinLib \texttt{.stp} file and prints one CSV row; \texttt{--tree=FILE} writes the tree's edges.
The README in the archive lists all options.

\section{Per-instance results}\label{app:tables}
Costs in bold equal the optimum. A superscript t marks a branch-and-bound run stopped by the 300\,s limit; its cost is the best tree
found. A dash in the DP columns means \(k>20\), which the DP does not attempt. Times are wall-clock seconds for one run.
\subsection*{Set B}
\begingroup\scriptsize\setlength{\tabcolsep}{3.2pt}
\begin{longtable}{@{}lrrrrrrrrrrrrrr@{}}\toprule
Inst. & \(n\) & \(m\) & \(k\) & OPT & \multicolumn{2}{c}{G1} & \multicolumn{2}{c}{G2} & \multicolumn{2}{c}{G3} & \multicolumn{2}{c}{DP} & \multicolumn{2}{c}{B\&B}\\
 & & & & & cost & s & cost & s & cost & s & cost & s & cost & s\\\midrule\endhead
b01 & 50 & 63 & 9 & 82 & \textbf{82} & 0.000 & \textbf{82} & 0.000 & \textbf{82} & 0.000 & \textbf{82} & 0.003 & \textbf{82} & 0.000\\
b02 & 50 & 63 & 13 & 83 & 88 & 0.000 & \textbf{83} & 0.000 & \textbf{83} & 0.000 & \textbf{83} & 0.100 & \textbf{83} & 0.000\\
b03 & 50 & 63 & 25 & 138 & \textbf{138} & 0.000 & \textbf{138} & 0.002 & \textbf{138} & 0.000 & -- & -- & \textbf{138} & 0.001\\
b04 & 50 & 100 & 9 & 59 & 62 & 0.000 & \textbf{59} & 0.002 & \textbf{59} & 0.000 & \textbf{59} & 0.002 & \textbf{59} & 0.003\\
b05 & 50 & 100 & 13 & 61 & \textbf{61} & 0.000 & \textbf{61} & 0.001 & \textbf{61} & 0.000 & \textbf{61} & 0.128 & \textbf{61} & 0.002\\
b06 & 50 & 100 & 25 & 122 & 127 & 0.002 & 124 & 0.006 & 124 & 0.000 & -- & -- & \textbf{122} & 0.010\\
b07 & 75 & 94 & 13 & 111 & \textbf{111} & 0.000 & \textbf{111} & 0.003 & \textbf{111} & 0.004 & \textbf{111} & 0.213 & \textbf{111} & 0.005\\
b08 & 75 & 94 & 19 & 104 & \textbf{104} & 0.000 & \textbf{104} & 0.001 & \textbf{104} & 0.000 & \textbf{104} & 27.2 & \textbf{104} & 0.001\\
b09 & 75 & 94 & 38 & 220 & 221 & 0.000 & \textbf{220} & 0.021 & \textbf{220} & 0.001 & -- & -- & \textbf{220} & 0.007\\
b10 & 75 & 150 & 13 & 86 & 91 & 0.000 & \textbf{86} & 0.001 & 87 & 0.000 & \textbf{86} & 0.145 & \textbf{86} & 0.003\\
b11 & 75 & 150 & 19 & 88 & 92 & 0.000 & \textbf{88} & 0.002 & \textbf{88} & 0.001 & \textbf{88} & 26.6 & \textbf{88} & 0.006\\
b12 & 75 & 150 & 38 & 174 & \textbf{174} & 0.002 & \textbf{174} & 0.012 & \textbf{174} & 0.003 & -- & -- & \textbf{174} & 0.049\\
b13 & 100 & 125 & 17 & 165 & 175 & 0.001 & \textbf{165} & 0.004 & 170 & 0.000 & \textbf{165} & 5.225 & \textbf{165} & 0.002\\
b14 & 100 & 125 & 25 & 235 & 237 & 0.005 & \textbf{235} & 0.017 & 236 & 0.007 & -- & -- & \textbf{235} & 0.003\\
b15 & 100 & 125 & 50 & 318 & \textbf{318} & 0.001 & \textbf{318} & 0.048 & \textbf{318} & 0.001 & -- & -- & \textbf{318} & 0.012\\
b16 & 100 & 200 & 17 & 127 & 135 & 0.000 & \textbf{127} & 0.007 & 133 & 0.001 & \textbf{127} & 4.922 & \textbf{127} & 0.024\\
b17 & 100 & 200 & 25 & 131 & 134 & 0.004 & 132 & 0.017 & 132 & 0.004 & -- & -- & \textbf{131} & 0.020\\
b18 & 100 & 200 & 50 & 218 & 221 & 0.002 & 220 & 0.061 & \textbf{218} & 0.002 & -- & -- & \textbf{218} & 0.037\\
\bottomrule\end{longtable}\endgroup
\subsection*{Set C}
\begingroup\scriptsize\setlength{\tabcolsep}{3.2pt}
\begin{longtable}{@{}lrrrrrrrrrrrrrr@{}}\toprule
Inst. & \(n\) & \(m\) & \(k\) & OPT & \multicolumn{2}{c}{G1} & \multicolumn{2}{c}{G2} & \multicolumn{2}{c}{G3} & \multicolumn{2}{c}{DP} & \multicolumn{2}{c}{B\&B}\\
 & & & & & cost & s & cost & s & cost & s & cost & s & cost & s\\\midrule\endhead
c01 & 500 & 625 & 5 & 85 & 88 & 0.001 & \textbf{85} & 0.002 & 87 & 0.002 & \textbf{85} & 0.010 & \textbf{85} & 0.030\\
c02 & 500 & 625 & 10 & 144 & \textbf{144} & 0.001 & \textbf{144} & 0.003 & \textbf{144} & 0.003 & \textbf{144} & 0.160 & \textbf{144} & 0.018\\
c03 & 500 & 625 & 83 & 754 & 769 & 0.020 & 761 & 0.315 & 761 & 0.052 & -- & -- & \textbf{754} & 0.386\\
c04 & 500 & 625 & 125 & 1079 & 1105 & 0.015 & 1088 & 0.902 & 1100 & 0.119 & -- & -- & \textbf{1079} & 0.604\\
c05 & 500 & 625 & 250 & 1579 & 1585 & 0.036 & 1581 & 3.505 & 1581 & 0.705 & -- & -- & \textbf{1579} & 1.115\\
c06 & 500 & 1000 & 5 & 55 & 60 & 0.001 & \textbf{55} & 0.003 & \textbf{55} & 0.001 & \textbf{55} & 0.005 & \textbf{55} & 0.223\\
c07 & 500 & 1000 & 10 & 102 & 115 & 0.003 & \textbf{102} & 0.011 & 107 & 0.003 & \textbf{102} & 0.167 & \textbf{102} & 0.256\\
c08 & 500 & 1000 & 83 & 509 & 526 & 0.022 & 512 & 0.953 & 514 & 0.086 & -- & -- & \textbf{509} & 0.901\\
c09 & 500 & 1000 & 125 & 707 & 721 & 0.026 & 714 & 0.990 & 713 & 0.106 & -- & -- & \textbf{707} & 4.717\\
c10 & 500 & 1000 & 250 & 1093 & 1106 & 0.039 & 1095 & 4.338 & 1098 & 0.833 & -- & -- & \textbf{1093} & 3.277\\
c11 & 500 & 2500 & 5 & 32 & 37 & 0.002 & 34 & 0.004 & 33 & 0.002 & \textbf{32} & 0.005 & \textbf{32} & 0.506\\
c12 & 500 & 2500 & 10 & 46 & 49 & 0.003 & \textbf{46} & 0.028 & 48 & 0.003 & \textbf{46} & 0.159 & \textbf{46} & 0.419\\
c13 & 500 & 2500 & 83 & 258 & 268 & 0.029 & 260 & 1.049 & 264 & 0.041 & -- & -- & \textbf{258} & 9.081\\
c14 & 500 & 2500 & 125 & 323 & 337 & 0.038 & 328 & 2.465 & 331 & 0.175 & -- & -- & \textbf{323} & 3.840\\
c15 & 500 & 2500 & 250 & 556 & 561 & 0.094 & 557 & 4.952 & 559 & 0.574 & -- & -- & \textbf{556} & 4.400\\
c16 & 500 & 12500 & 5 & 11 & 12 & 0.002 & 12 & 0.005 & 12 & 0.003 & \textbf{11} & 0.007 & \textbf{11} & 0.451\\
c17 & 500 & 12500 & 10 & 18 & 20 & 0.005 & 19 & 0.017 & 19 & 0.005 & \textbf{18} & 0.238 & \textbf{18} & 0.465\\
c18 & 500 & 12500 & 83 & 113 & 128 & 0.037 & 120 & 0.827 & 119 & 0.050 & -- & -- & 117\textsuperscript{t} & 300.0\\
c19 & 500 & 12500 & 125 & 146 & 154 & 0.026 & 155 & 2.756 & 148 & 0.084 & -- & -- & \textbf{146} & 7.929\\
c20 & 500 & 12500 & 250 & 267 & 268 & 0.121 & 268 & 10.9 & \textbf{267} & 0.230 & -- & -- & \textbf{267} & 4.765\\
\bottomrule\end{longtable}\endgroup
\subsection*{Set P6E}
\begingroup\scriptsize\setlength{\tabcolsep}{3.2pt}
\begin{longtable}{@{}lrrrrrrrrrrrrrr@{}}\toprule
Inst. & \(n\) & \(m\) & \(k\) & OPT & \multicolumn{2}{c}{G1} & \multicolumn{2}{c}{G2} & \multicolumn{2}{c}{G3} & \multicolumn{2}{c}{DP} & \multicolumn{2}{c}{B\&B}\\
 & & & & & cost & s & cost & s & cost & s & cost & s & cost & s\\\midrule\endhead
p619 & 100 & 180 & 5 & 7485 & \textbf{7485} & 0.000 & \textbf{7485} & 0.000 & \textbf{7485} & 0.000 & \textbf{7485} & 0.001 & \textbf{7485} & 0.007\\
p620 & 100 & 180 & 5 & 8746 & 8853 & 0.000 & \textbf{8746} & 0.000 & \textbf{8746} & 0.000 & \textbf{8746} & 0.001 & \textbf{8746} & 0.006\\
p621 & 100 & 180 & 5 & 8688 & 8791 & 0.000 & \textbf{8688} & 0.000 & \textbf{8688} & 0.000 & \textbf{8688} & 0.001 & \textbf{8688} & 0.007\\
p622 & 100 & 180 & 10 & 15972 & 17403 & 0.000 & \textbf{15972} & 0.001 & \textbf{15972} & 0.000 & \textbf{15972} & 0.022 & \textbf{15972} & 0.009\\
p623 & 100 & 180 & 10 & 19496 & 20667 & 0.000 & \textbf{19496} & 0.002 & \textbf{19496} & 0.000 & \textbf{19496} & 0.023 & \textbf{19496} & 0.011\\
p624 & 100 & 180 & 20 & 20246 & \textbf{20246} & 0.001 & \textbf{20246} & 0.005 & \textbf{20246} & 0.001 & \textbf{20246} & 150.0 & \textbf{20246} & 0.045\\
p625 & 100 & 180 & 20 & 23078 & 24193 & 0.001 & \textbf{23078} & 0.012 & 23098 & 0.001 & \textbf{23078} & 155.1 & \textbf{23078} & 0.039\\
p626 & 100 & 180 & 20 & 22346 & 23000 & 0.003 & \textbf{22346} & 0.014 & \textbf{22346} & 0.002 & \textbf{22346} & 140.8 & \textbf{22346} & 0.043\\
p627 & 100 & 180 & 50 & 40647 & 40748 & 0.003 & \textbf{40647} & 0.073 & \textbf{40647} & 0.001 & -- & -- & \textbf{40647} & 0.088\\
p628 & 100 & 180 & 50 & 40008 & 40607 & 0.001 & \textbf{40008} & 0.041 & \textbf{40008} & 0.002 & -- & -- & \textbf{40008} & 0.140\\
p629 & 100 & 180 & 50 & 43287 & 43746 & 0.001 & 43335 & 0.035 & \textbf{43287} & 0.002 & -- & -- & \textbf{43287} & 0.107\\
p630 & 200 & 370 & 10 & 26125 & 26761 & 0.001 & 26416 & 0.002 & \textbf{26125} & 0.001 & \textbf{26125} & 0.047 & \textbf{26125} & 0.033\\
p631 & 200 & 370 & 20 & 39067 & 42103 & 0.001 & 40429 & 0.011 & 39565 & 0.002 & -- & -- & \textbf{39067} & 2.407\\
p632 & 200 & 370 & 40 & 56217 & 57776 & 0.003 & 56562 & 0.052 & \textbf{56217} & 0.003 & -- & -- & \textbf{56217} & 13.1\\
p633 & 200 & 370 & 100 & 86268 & 86877 & 0.007 & 86297 & 0.336 & \textbf{86268} & 0.012 & -- & -- & \textbf{86268} & 0.523\\
\bottomrule\end{longtable}\endgroup
\subsection*{Set Adversarial}
\begingroup\scriptsize\setlength{\tabcolsep}{3.2pt}
\begin{longtable}{@{}lrrrrrrrrrrrrrr@{}}\toprule
Inst. & \(n\) & \(m\) & \(k\) & OPT & \multicolumn{2}{c}{G1} & \multicolumn{2}{c}{G2} & \multicolumn{2}{c}{G3} & \multicolumn{2}{c}{DP} & \multicolumn{2}{c}{B\&B}\\
 & & & & & cost & s & cost & s & cost & s & cost & s & cost & s\\\midrule\endhead
adv\_g1 & 7 & 8 & 4 & 30 & 42 & 0.000 & \textbf{30} & 0.000 & \textbf{30} & 0.000 & \textbf{30} & 0.000 & \textbf{30} & 0.000\\
adv\_g2 & 8 & 10 & 5 & 49 & \textbf{49} & 0.000 & 68 & 0.000 & \textbf{49} & 0.000 & \textbf{49} & 0.000 & \textbf{49} & 0.000\\
adv\_g3 & 9 & 9 & 5 & 48 & 58 & 0.000 & \textbf{48} & 0.000 & 58 & 0.000 & \textbf{48} & 0.000 & \textbf{48} & 0.000\\
star\_l03 & 4 & 6 & 3 & 30 & 38 & 0.000 & 38 & 0.000 & \textbf{30} & 0.000 & \textbf{30} & 0.000 & \textbf{30} & 0.000\\
star\_l04 & 5 & 10 & 4 & 40 & 57 & 0.000 & 57 & 0.000 & \textbf{40} & 0.000 & \textbf{40} & 0.000 & \textbf{40} & 0.000\\
star\_l05 & 6 & 15 & 5 & 50 & 76 & 0.000 & 76 & 0.000 & \textbf{50} & 0.000 & \textbf{50} & 0.000 & \textbf{50} & 0.000\\
star\_l06 & 7 & 21 & 6 & 60 & 95 & 0.000 & 95 & 0.000 & \textbf{60} & 0.000 & \textbf{60} & 0.000 & \textbf{60} & 0.000\\
star\_l07 & 8 & 28 & 7 & 70 & 114 & 0.000 & 114 & 0.000 & \textbf{70} & 0.000 & \textbf{70} & 0.000 & \textbf{70} & 0.000\\
star\_l08 & 9 & 36 & 8 & 80 & 133 & 0.000 & 133 & 0.000 & \textbf{80} & 0.000 & \textbf{80} & 0.000 & \textbf{80} & 0.000\\
star\_l09 & 10 & 45 & 9 & 90 & 152 & 0.000 & 152 & 0.001 & \textbf{90} & 0.000 & \textbf{90} & 0.001 & \textbf{90} & 0.000\\
star\_l10 & 11 & 55 & 10 & 100 & 171 & 0.000 & 171 & 0.000 & \textbf{100} & 0.000 & \textbf{100} & 0.001 & \textbf{100} & 0.001\\
star\_l11 & 12 & 66 & 11 & 110 & 190 & 0.000 & 190 & 0.000 & \textbf{110} & 0.000 & \textbf{110} & 0.005 & \textbf{110} & 0.001\\
star\_l12 & 13 & 78 & 12 & 120 & 209 & 0.000 & 209 & 0.000 & \textbf{120} & 0.000 & \textbf{120} & 0.007 & \textbf{120} & 0.001\\
\bottomrule\end{longtable}\endgroup
\subsection*{Set Grid}
\begingroup\scriptsize\setlength{\tabcolsep}{3.2pt}
\begin{longtable}{@{}lrrrrrrrrrrrrrr@{}}\toprule
Inst. & \(n\) & \(m\) & \(k\) & OPT & \multicolumn{2}{c}{G1} & \multicolumn{2}{c}{G2} & \multicolumn{2}{c}{G3} & \multicolumn{2}{c}{DP} & \multicolumn{2}{c}{B\&B}\\
 & & & & & cost & s & cost & s & cost & s & cost & s & cost & s\\\midrule\endhead
grid\_k04\_s0 & 100 & 180 & 4 & 89 & 94 & 0.000 & 90 & 0.000 & \textbf{89} & 0.000 & \textbf{89} & 0.000 & \textbf{89} & 0.023\\
grid\_k04\_s1 & 100 & 180 & 4 & 33 & \textbf{33} & 0.000 & \textbf{33} & 0.000 & \textbf{33} & 0.000 & \textbf{33} & 0.001 & \textbf{33} & 0.014\\
grid\_k04\_s2 & 100 & 180 & 4 & 34 & \textbf{34} & 0.000 & \textbf{34} & 0.000 & \textbf{34} & 0.000 & \textbf{34} & 0.000 & \textbf{34} & 0.010\\
grid\_k06\_s0 & 100 & 180 & 6 & 96 & 103 & 0.000 & \textbf{96} & 0.001 & \textbf{96} & 0.000 & \textbf{96} & 0.001 & \textbf{96} & 0.005\\
grid\_k06\_s1 & 100 & 180 & 6 & 77 & \textbf{77} & 0.000 & \textbf{77} & 0.001 & \textbf{77} & 0.000 & \textbf{77} & 0.001 & \textbf{77} & 0.013\\
grid\_k06\_s2 & 100 & 180 & 6 & 61 & \textbf{61} & 0.000 & 62 & 0.000 & \textbf{61} & 0.000 & \textbf{61} & 0.001 & \textbf{61} & 0.014\\
grid\_k08\_s0 & 100 & 180 & 8 & 117 & 118 & 0.000 & 118 & 0.000 & \textbf{117} & 0.000 & \textbf{117} & 0.002 & \textbf{117} & 0.005\\
grid\_k08\_s1 & 100 & 180 & 8 & 94 & 98 & 0.000 & \textbf{94} & 0.000 & 96 & 0.000 & \textbf{94} & 0.002 & \textbf{94} & 0.013\\
grid\_k08\_s2 & 100 & 180 & 8 & 85 & 93 & 0.000 & 86 & 0.001 & \textbf{85} & 0.000 & \textbf{85} & 0.005 & \textbf{85} & 0.005\\
grid\_k10\_s0 & 100 & 180 & 10 & 107 & 116 & 0.000 & \textbf{107} & 0.002 & 108 & 0.000 & \textbf{107} & 0.022 & \textbf{107} & 0.016\\
grid\_k10\_s1 & 100 & 180 & 10 & 113 & 115 & 0.000 & \textbf{113} & 0.002 & 114 & 0.000 & \textbf{113} & 0.022 & \textbf{113} & 0.017\\
grid\_k10\_s2 & 100 & 180 & 10 & 94 & \textbf{94} & 0.000 & \textbf{94} & 0.001 & \textbf{94} & 0.000 & \textbf{94} & 0.022 & \textbf{94} & 0.010\\
grid\_k12\_s0 & 100 & 180 & 12 & 98 & 101 & 0.000 & \textbf{98} & 0.002 & \textbf{98} & 0.000 & \textbf{98} & 0.095 & \textbf{98} & 0.015\\
grid\_k12\_s1 & 100 & 180 & 12 & 106 & 113 & 0.001 & \textbf{106} & 0.002 & 107 & 0.001 & \textbf{106} & 0.097 & \textbf{106} & 0.014\\
grid\_k12\_s2 & 100 & 180 & 12 & 121 & 122 & 0.000 & 122 & 0.002 & \textbf{121} & 0.001 & \textbf{121} & 0.097 & \textbf{121} & 0.020\\
grid\_k14\_s0 & 100 & 180 & 14 & 126 & 130 & 0.001 & 129 & 0.003 & 129 & 0.001 & \textbf{126} & 0.493 & \textbf{126} & 0.017\\
grid\_k14\_s1 & 100 & 180 & 14 & 147 & 156 & 0.001 & \textbf{147} & 0.003 & 151 & 0.001 & \textbf{147} & 0.462 & \textbf{147} & 0.014\\
grid\_k14\_s2 & 100 & 180 & 14 & 111 & 118 & 0.001 & 112 & 0.002 & 113 & 0.001 & \textbf{111} & 0.440 & \textbf{111} & 0.014\\
grid\_k16\_s0 & 100 & 180 & 16 & 121 & 125 & 0.001 & \textbf{121} & 0.004 & 123 & 0.001 & \textbf{121} & 1.363 & \textbf{121} & 0.008\\
grid\_k16\_s1 & 100 & 180 & 16 & 123 & 124 & 0.001 & \textbf{123} & 0.004 & \textbf{123} & 0.001 & \textbf{123} & 3.380 & \textbf{123} & 0.020\\
grid\_k16\_s2 & 100 & 180 & 16 & 109 & \textbf{109} & 0.001 & \textbf{109} & 0.003 & \textbf{109} & 0.001 & \textbf{109} & 3.531 & \textbf{109} & 0.013\\
grid\_k18\_s0 & 100 & 180 & 18 & 112 & 115 & 0.001 & 115 & 0.005 & \textbf{112} & 0.003 & \textbf{112} & 16.1 & \textbf{112} & 0.022\\
grid\_k18\_s1 & 100 & 180 & 18 & 119 & 121 & 0.001 & 121 & 0.004 & \textbf{119} & 0.001 & \textbf{119} & 14.3 & \textbf{119} & 0.021\\
grid\_k18\_s2 & 100 & 180 & 18 & 151 & 154 & 0.001 & 153 & 0.005 & 152 & 0.001 & \textbf{151} & 16.9 & \textbf{151} & 0.040\\
grid\_k20\_s0 & 100 & 180 & 20 & 166 & 169 & 0.001 & \textbf{166} & 0.012 & 168 & 0.001 & \textbf{166} & 124.7 & \textbf{166} & 0.020\\
grid\_k20\_s1 & 100 & 180 & 20 & 139 & 140 & 0.001 & \textbf{139} & 0.005 & \textbf{139} & 0.001 & \textbf{139} & 113.7 & \textbf{139} & 0.017\\
grid\_k20\_s2 & 100 & 180 & 20 & 151 & 153 & 0.003 & \textbf{151} & 0.003 & \textbf{151} & 0.001 & \textbf{151} & 107.7 & \textbf{151} & 0.024\\
\bottomrule\end{longtable}\endgroup
\subsection*{Set Random}
\begingroup\scriptsize\setlength{\tabcolsep}{3.2pt}
\begin{longtable}{@{}lrrrrrrrrrrrrrr@{}}\toprule
Inst. & \(n\) & \(m\) & \(k\) & OPT & \multicolumn{2}{c}{G1} & \multicolumn{2}{c}{G2} & \multicolumn{2}{c}{G3} & \multicolumn{2}{c}{DP} & \multicolumn{2}{c}{B\&B}\\
 & & & & & cost & s & cost & s & cost & s & cost & s & cost & s\\\midrule\endhead
rand\_n0100 & 100 & 200 & 10 & 689 & 699 & 0.000 & 697 & 0.002 & 691 & 0.001 & \textbf{689} & 0.024 & \textbf{689} & 0.011\\
rand\_n0200 & 200 & 400 & 10 & 827 & 945 & 0.001 & \textbf{827} & 0.004 & 864 & 0.001 & \textbf{827} & 0.058 & \textbf{827} & 0.033\\
rand\_n0400 & 400 & 800 & 10 & 1226 & 1370 & 0.002 & 1244 & 0.011 & 1291 & 0.002 & \textbf{1226} & 0.177 & \textbf{1226} & 0.176\\
rand\_n0800 & 800 & 1600 & 10 & 2005 & 2194 & 0.003 & 2046 & 0.018 & 2046 & 0.003 & \textbf{2005} & 0.246 & \textbf{2005} & 0.670\\
rand\_n1600 & 1600 & 3200 & 10 & 2652 & 2909 & 0.007 & 2691 & 0.034 & 2706 & 0.007 & \textbf{2652} & 0.525 & \textbf{2652} & 2.086\\
rand\_n3200 & 3200 & 6400 & 10 & 4556 & 4606 & 0.014 & 4557 & 0.065 & 4574 & 0.014 & \textbf{4556} & 1.116 & \textbf{4556} & 8.735\\
\bottomrule\end{longtable}\endgroup
\subsection*{Set I080 (100 instances, grouped)}
\begin{table}[H]\centering\small\begin{tabular}{@{}rrrrrrrrrrrr@{}}\toprule
\(k\) & \# & \multicolumn{2}{c}{G1} & \multicolumn{2}{c}{G2} & \multicolumn{2}{c}{G3} & \multicolumn{2}{c}{DP} & \multicolumn{2}{c}{B\&B}\\
 & & gap \% & opt & gap \% & opt & gap \% & opt & solved & max s & solved & max s\\\midrule
6 & 25 & 14.83 & 0 & 7.76 & 5 & 7.53 & 2 & 25 & 0.01 & 25 & 0.03\\
8 & 25 & 18.28 & 1 & 8.42 & 5 & 6.44 & 1 & 25 & 0.01 & 25 & 0.04\\
16 & 25 & 19.66 & 0 & 12.83 & 1 & 9.27 & 0 & 25 & 4.40 & 25 & 0.29\\
20 & 25 & 21.27 & 0 & 13.93 & 2 & 9.32 & 0 & 25 & 159.31 & 25 & 0.53\\
\bottomrule\end{tabular}\end{table}



\end{document}
